\providecommand{\toplevelprefix}{../..}  %
\documentclass[../../book-main_es.tex]{subfiles}

\begin{document}

\chapter{Entropía, difusión, eliminación de ruido y
codificación con pérdida}\label{app:entropy}

\begin{quote}
    ``{\em El aumento del desorden o la entropía con el paso del tiempo es un
        ejemplo de lo que se denomina «flecha del tiempo», algo que
    distingue el pasado del futuro, dando una dirección al tiempo.}''

    $~$\hfill -- Breve historia del tiempo, Stephen Hawking
\end{quote}
\vspace{5mm}

En este apéndice ofrecemos demostraciones de varios resultados, mencionados en el
\Cref{ch:denoising}, relacionados con la entropía diferencial, cómo
evoluciona en procesos de difusión y sus conexiones con la codificación con pérdida. Haremos
la siguiente supuesto moderado sobre la variable aleatoria que representa
la fuente de datos, denotada por \(\vx\).

\begin{assumption}\label{assumption:entropy_x_compact_support}
    \(\vx\) se define en un conjunto compacto \(\cS \subseteq \R^{D}\) de
    radio no mayor que \(R\), es decir, \(R \doteq \sup_{\vxi \in \cS}\norm{\vxi}_{2}\).
\end{assumption}

En particular, dado que los conjuntos compactos en el espacio euclídeo son acotados,
se cumple que \(R < \infty\). Utilizaremos de manera sistemática la notación
\(B_{r}(\vxi) \doteq \{\vu \in \R^{D} \colon \norm{\vxi - \vu}_{2}
\leq r\}\) para denotar la bola euclidiana de radio \(r\) centrada en
\(\vxi\). En este sentido, el \Cref{assumption:entropy_x_compact_support}
tiene \(\cS \subseteq B_{R}(\vzero)\).

Obsérvese que este supuesto se cumple para (casi) todas las variables que nos interesan
en la práctica, ya que (a menudo) viene impuesta por un paso de normalización
durante el preprocesamiento de los datos.

\section{Entropía diferencial de distribuciones
de baja dimensión}\label{sec:low_dim_entropy}

En este breve apéndice analizamos la entropía diferencial de
las distribuciones de baja dimensión. Por definición, la entropía
diferencial de una variable aleatoria \(\vx\) que no posee densidad
es \(-\infty\); esto incluye a todas las variables aleatorias definidas en
conjuntos de baja dimensión. El objetivo de esta sección es analizar por qué
este es un valor «moralmente correcto».

De hecho, sea \(\vx\) cualquier variable aleatoria tal que
se cumpla la hipótesis del \Cref{assumption:entropy_x_compact_support}, es decir, el soporte
\(\cS\) de \(\vx\) tiene volumen \(0\).\footnote{Formalmente, esto significa
que \(\cS\) es boreliana medible con medida boreliana \(0\).} Consideraremos
el caso en que \(\vx\) sea uniforme en \(\cS\).\footnote{Digamos,
con respecto a la medida de Hausdorff en \(\cS\).} Nuestro objetivo es calcular \(h(\vx)\).

En este caso, \(\vx\) no tendría densidad; en el caso hipotético
en el que no supiéramos que \(h(\vx) = -\infty\), no podríamos
definirla directamente utilizando la definición estándar de
entropía diferencial. En cambio, como en el resto del análisis y la teoría de la información,
sería razonable considerar el límite de las entropías de 
aproximaciones sucesivamente mejores \(\vx_{\eps}\) de \(\vx\) que
tienen densidades, es decir,
\begin{equation}
    h(\vx)\ \text{``=''}\ \lim_{\eps \searrow 0}h(\vx_{\eps}).
\end{equation}

Con este fin, la idea básica consiste en tomar un entorno tubular de espesor \(\eps\) de
\(\cS\), denotado por ejemplo como \(\cS_{\eps}\), definido como
\begin{equation}
    \cS_{\eps} = \bigcup_{\vxi \in \cS}B_{\eps}(\vxi)
\end{equation}
y que se muestra en la \Cref{fig:entropy_eps_thickening}.
\begin{figure}[th]
    \centering
    \begin{tikzpicture}
        \def\radius{0.2cm} %
        \def\curve{(0,0) .. controls (1,1.5) and (3,-0.5) .. (4,1)} %

        \path[decoration={markings, mark=between positions 0 and 1
                step 0.02 with {
                    \fill[red!30, opacity=0.5, draw=none] (0,0)
                    circle (\radius);
        }}, postaction={decorate}] \curve;

        \draw[blue, thick] \curve node[right, blue] {$\cS$};

        \coordinate (p_on_curve) at (2, 0.5); %
        \fill[red!50, opacity=0.7] (p_on_curve) circle (\radius);
        \fill[black] (p_on_curve) circle (1pt) node[below left] {$x$};

        \node[red, below right] at (3.5, -0.2) {$\cS_{\eps} =
        \bigcup_{\vxi \in \cS} B_{\eps}(\vxi)$};

    \end{tikzpicture}
    \caption{Ilustración del entorno tubular de espesor \(\eps\) \(\cS_\eps\) de
    una curva \(\cS \subseteq \R^{2}\).}
    \label{fig:entropy_eps_thickening}
\end{figure}
Trabajaremos con variables aleatorias cuyo soporte sea \(\cS_{\eps}\),
que es de dimensión completa, y tomaremos el límite cuando \(\eps \to 0\).
De hecho, definamos \(\vx_{\eps} \sim \dUnif(\cS_{\eps})\). Dado que
\(\cS_{\eps}\) tiene medida de Lebesgue positiva, \(\vx_{\eps}\) tiene una densidad
\(p_{\eps}\) igual a
\begin{equation}
    p_{\eps}(\vxi) = \indvar(\vxi \in \cS_{\eps}) \cdot
    \frac{1}{\volume(\cS_{\eps})}.
\end{equation}
Al calcular la entropía de \(\vx_{\eps}\) utilizando la convención de que \(0
\log 0 = 0\), se cumple que
\begin{align}
    h(\vx_{\eps})
    &= -\int_{\R^{D}}p_{\eps}(\vxi)\log p_{\eps}(\vxi)\odif{\vxi} \\
    &= -\int_{\cS_{\eps}}\frac{1}{\volume(\cS_{\eps})}
    \log\rp{\frac{1}{\volume(\cS_{\eps})}}\odif{\vxi} \\
    &=
    \frac{\log(\volume(\cS_{\eps}))}{\volume(\cS_{\eps})}\int_{\cS_{\eps}}\odif{\vxi}
    \\
    &= \log(\volume(\cS_{\eps})).
\end{align}
Dado que \(\cS\) es compacto, \(\volume(\cS_{\eps})\) es finito y tiende
a \(0\) cuando \(\eps \searrow 0\). Por lo tanto
\begin{equation}
    h(\vx) = \lim_{\eps \searrow 0}h(\vx_{\eps}) = \lim_{\eps
    \searrow 0}\log(\volume(\cS_{\eps})) = -\infty,
\end{equation}
tal como se desea.

El cálculo anterior es, en realidad, una consecuencia de un conjunto de resultados mucho más famosos
y reconocidos sobre la entropía máxima posible de
\(\vx\) sujeta a ciertas restricciones en la distribución de
\(\vx\). Sería una omisión por nuestra parte no incluir aquí dichos resultados; las
demostraciones se encuentran, por ejemplo, en el capítulo 2 de \cite{poliyanski2024information},
por ejemplo.
\begin{theorem}\label{thm:max_entropy}
    Sea \(\vx\) una variable aleatoria en \(\R^{D}\).
    \begin{enumerate}
        \item Si \(\vx\) tiene soporte en un conjunto compacto \(\cS
            \subseteq \R^{D}\) (es decir,
            \Cref{assumption:entropy_x_compact_support}), entonces
            \begin{equation}
                h(\vx) \leq h(\dUnif(\cS)) = \log \volume(\cS).
            \end{equation}
        \item Si \(\vx\) tiene matriz de covarianza finita tal que, para una matriz PSD
            \(\vSigma \in \PSD(D)\), se cumple que \(\Cov(\vx)
            \preceq \vSigma\) (con respecto al orden PSD, es decir,
            \(\vSigma - \Cov(\vx)\) es PSD), entonces
            \begin{equation}
                h(\vx) \leq h(\dNorm(\vzero, \vSigma)) =
                \frac{1}{2}\log((2\pi e)^{D}\det\vSigma).
            \end{equation}
        \item Si \(\vx\) tiene un momento de segundo orden finito tal que, para una
            constante \(a \geq 0\), se cumple que \(\Ex\norm{\vx}_{2}^{2}
            \leq a\), entonces
            \begin{equation}
                h(\vx) \leq h\rp{\dNorm\rp{\vzero, \frac{a}{D}\vI}} =
                \frac{D}{2}\log\frac{2\pi e a}{D}.
            \end{equation}
    \end{enumerate}
\end{theorem}

\section{Procesos de difusión y eliminación de ruido}\label{sec:entropy_diffusion}

En el cuerpo principal (\Cref{ch:denoising}), consideramos una variable aleatoria
\(\vx\) y un proceso estocástico definido por
\eqref{eq:additive_gaussian_noise_model}, es decir,
\begin{equation}\label{eq:app_additive_gaussian_noise_model}
    \vx_{t} = \vx + t\vg, \qquad  \forall t \in [0, T]
\end{equation}
donde \(\vg \sim \dNorm(\vzero, \vI)\) independientemente de \(\vx\).

La estructura de esta sección es la siguiente. En el
\Cref{sub:diffusion_entropy_increases} presentamos un teorema formal
y una demostración rigurosa que muestra que, bajo la
\Cref{eq:app_additive_gaussian_noise_model}, la entropía aumenta,
es decir, \(\odv*{h(\vx_{t})}{t} > 0\). En el
\Cref{sub:denoising_entropy_decreases} presentamos un teorema formal
y una demostración rigurosa que muestra que, bajo la
\Cref{eq:app_additive_gaussian_noise_model}, la entropía disminuye
durante el proceso de eliminación de ruido, es decir, \(h(\Ex[\vx_{s} \given \vx_{t}]) <
h(\vx_{t})\) para todo \(s < t\). En el
\Cref{sub:app_diffusion_intermediate_results} ofrecemos pruebas de
los lemas técnicos necesarios para establecer las afirmaciones de las
subsecciones anteriores.

Antes de empezar, vamos a definir algunas notaciones clave. En primer lugar, sea
\(\phi_{t}\) la densidad de \(\dNorm(\vzero, t^{2}\vI)\), es decir,
\begin{equation}\label{eq:gaussian_noise_time_t}
    \phi_{t}(\vxi) \doteq
    \frac{1}{(2\pi)^{D/2}t^{D}}\exp\rp{-\frac{\norm{\vxi}_{2}^{2}}{2t^{2}}}.
\end{equation}
A continuación, dado que \(\vx_{t}\) se define en todo el espacio \(\R^{D}\), tiene una
\textit{densidad}, que denotamos como \(p_{t}\) (al igual que en el cuerpo principal). Un
cálculo rápido muestra que
\begin{equation}\label{eq:p_t_representation}
    p_{t}(\vxi) = \Ex[\phi_{t}(\vxi - \vx)],
\end{equation}
y a partir de esta representación es posible deducir (es decir, a partir de
\Cref{prop:diff_convolution}) que \(p_{t}\) es suave (es decir,
infinitamente diferenciable) en \(\vxi\), también suave en \(t\), y
positiva en todas partes. Este hecho resulta algo notable a primera vista:
incluso para una variable aleatoria completamente irregular \(\vx\) (por ejemplo, una
variable aleatoria de Bernoulli, que no posee densidad), su
suavizado gaussiano admite una densidad para todo \(t
> 0\) (arbitrariamente pequeño). La demostración se deja como ejercicio para los lectores versados en
análisis matemático.

Sin embargo, también debemos añadir un supuesto sobre la
\textit{suavidad} de la distribución de \(\vx\), lo que
eliminará algunos problemas técnicos que surgen en torno a \(t = 0\) con
distribuciones de baja dimensión.\footnote{Ya que entonces diversas magnitudes
    se vuelven muy irregulares y manejarlas requeriría
un análisis adicional significativo.} A pesar de esto, esperamos que nuestros
resultados se mantengan bajo supuestos más suaves con trabajo adicional. Por ahora,
supongamos que:
\begin{assumption}\label{assumption:entropy_x_density}
    \(\vx\) tiene una densidad dos veces continuamente diferenciable, denotada por \(p\).
\end{assumption}

\subsection{El proceso de difusión aumenta la entropía con el
paso del tiempo}\label{sub:diffusion_entropy_increases}

En esta sección ofrecemos una demostración del
\Cref{thm:diffusion_entropy_increases}. Por razones de conveniencia, la reformulamos
de la siguiente manera.

\begin{theorem}[La difusión aumenta
    la entropía]\label{thm:diffusion_entropy_increases}
    Sea \(\vx\) una variable aleatoria cualquiera tal que los
    \Cref{assumption:entropy_x_compact_support,assumption:entropy_x_density}
    se cumplan, y sea \((\vx_{t})_{t \in [0, T]}\) el proceso estocástico
    \eqref{eq:app_additive_gaussian_noise_model}. Entonces
    \begin{equation}\label{eq:diffusion_entropy_increases}
        h(\vx_{s}) < h(\vx_{t}), \qquad \forall s, t \colon 0 \leq s < t \leq T.
    \end{equation}
\end{theorem}
\begin{proof}
    Antes de comenzar, debemos preguntarnos: ¿en qué casos tiene sentido la desigualdad de
    \eqref{eq:diffusion_entropy_increases}? Demostraremos
    en el \Cref{lem:diffusion_entropy_exists} que, bajo nuestras
    suposiciones, la entropía diferencial está bien definida, nunca es
    \(+\infty\) y, para \(t > 0\), es finita, por lo que la desigualdad (estricta)
    en \eqref{eq:diffusion_entropy_increases} tiene sentido.

    Dejando a un lado la cuestión de la biendefinición, el quid de esta demostración consiste
    en demostrar que la densidad \(p_{t}\) de \(\vx_{t}\) satisface una
    ecuación diferencial parcial concreta, que es muy similar
    a la \textit{ecuación del calor}. La ecuación del calor es una famosa ecuación diferencial parcial
    que describe la difusión del calor en el espacio. Esto
    debería tener sentido de manera intuitiva y nos permite formarnos una imagen mental: a medida que
    el tiempo \(t\) aumenta, la probabilidad del vector original
    (quizás muy concentrado) \(\vx\) se dispersa por todo
    \(\R^{D}\), como el calor que se irradia desde una fuente en el vacío.

    Estas ecuaciones diferenciales parciales para \(p_{t}\), conocidas como \textit{ecuaciones de Fokker-Planck}
    para procesos estocásticos más generales, son herramientas muy
    poderosas, ya que nos permiten describir las derivadas temporales instantáneas
    de \(p_{t}\) en términos de las derivadas espaciales instantáneas
    de \(p_{t}\), y viceversa, proporcionando una
    descripción concisa de la regularidad y la dinámica de \(p_{t}\).
    Una vez que obtenemos la dinámica de \(p_{t}\), podemos utilizar el sistema
    para obtener otra que describa la dinámica de
    \(h(\vx_{t})\), que, al fin y al cabo, no es más que una funcional de la densidad \(p_{t}\).

    La descripción de la EDP implica un objeto matemático denominado
    el laplaciano \(\Delta\). Recordemos de la clase de cálculo multivariable
    que el laplaciano, al actuar sobre una función \(f \colon (0, T)
    \times \R^{D} \to \R\) diferenciable en el tiempo
    y dos veces diferenciable en el espacio, viene dado por
    \begin{equation}
        \Delta f_{t}(\vxi) = \tr(\nabla^{2}f_{t}(\vxi)) = \sum_{i =
        1}^{D}\pddv{f_{t}}{\xi_{i}}(\vxi).
    \end{equation}

    Es decir, al utilizar la representación integral de \(p_{t}\) y
    derivarla con respecto a la integral, podemos calcular las
    derivadas de \(p_{t}\) (lo cual hacemos en
    \Cref{prop:p_t_derivatives}) y observamos que \(p_{t}\) satisface
    la ecuación diferencial parcial de tipo calorífico
    \begin{equation}
        \pdv{p_{t}}{t}(\vxi) = t\Delta p_{t}(\vxi).
    \end{equation}
    A continuación, para hallar la dinámica de \(h(\vx_{t})\), podemos volver a utilizar
    \Cref{prop:dutis}, así como la ecuación diferencial parcial de tipo calorífico, para obtener
    \begin{align}
        \odv*{h(\vx_{t})}{t}
        &= -\odv*{\int_{\R^{D}}p_{t}(\vxi)\log p_{t}(\vxi)\odif{\vxi}}{t} \\
        &= -\int_{\R^{D}}\pdv*{\bs{p_{t}(\vxi)\log
        p_{t}(\vxi)}}{t}\odif{\vxi} \\
        &= -\int_{\R^{D}}\pdv{p_{t}}{t}(\vxi)[1 + \log
        p_{t}(\vxi)]\odif{\vxi} \\
        &= -t\int_{\R^{D}}\Delta p_{t}(\vxi)[1 + \log p_{t}(\vxi)]\odif{\vxi}.
    \end{align}
    Mediante un razonamiento ligeramente complejo basado en la integración por partes
    (\Cref{lem:diffusion_ibp}), obtenemos
    \begin{align}
        \odv*{h(\vx_{t})}{t}
        &= t\int_{\R^{D}}\ip{\nabla \log p_{t}(\vxi)}{\nabla
        p_{t}(\vxi)}\odif{\vxi} \\
        &= t\int_{\R^{D}}\frac{\norm{\nabla
        p_{t}(\vxi)}_{2}^{2}}{p_{t}(\vxi)}\odif{\vxi} \\
        &> 0
    \end{align}
    donde en la última línea se cumple una desigualdad estricta porque, para que
    no se cumpliera, \(\nabla p_{t}(\vxi)\) tendría que ser nulo casi
    en todas partes (es decir, en todas partes excepto, posiblemente, en un conjunto de volumen cero), 
    pero esto implicaría que \(p_{t}\) sería constante
    casi en todas partes, lo cual contradice el hecho de que \(p_{t}\) es una densidad.

    Para completar la demostración, basta con aplicar el teorema fundamental del cálculo
    \begin{equation}
        h(\vx_{t}) = h(\vx_{s}) +
        \int_{s}^{t}\odv*{h(\vx_{u})}{u}\odif{u} > h(\vx_{s}),
    \end{equation}
    lo cual demuestra la afirmación. (Obsérvese que esto no tiene sentido cuando
        \(h(\vx_{s}) = -\infty\), lo cual solo puede ocurrir cuando \(s = 0\)
        y \(h(\vx) = -\infty\), pero en este caso \(h(\vx_{t}) >
    -\infty\), por lo que la afirmación es verdadera de manera vacía de todos modos.)
\end{proof}

\subsection{El proceso de eliminación de ruido reduce la entropía con el
paso del tiempo}\label{sub:denoising_entropy_decreases}

Recordemos que en la \Cref{sub:intro_diffusion_denoising} partimos de la
variable aleatoria \(\vx_{T}\) y eliminamos el ruido de forma iterativa mediante
iteraciones de la forma
\begin{equation}\label{eq:app_denoising_iteration}
    \hat{\vx}_{s} \doteq \Ex[\vx_{s} \mid \vx_{t} = \hat{\vx}_{t}] =
    \frac{s}{t}\hat{\vx}_{t} + \bp{1 -
    \frac{s}{t}}\bar{\vx}^{\ast}(t, \hat{\vx}_{t}).
\end{equation}
para \(s, t \in \{t_{0}, t_{1}, \dots, t_{L}\}\) con \(s < t\) y
\(\vx_{T} = \hat{\vx}_{T}\). Queremos demostrar que \(h(\hat{\vx}_{s})
< h(\hat{\vx}_{t})\), demostrando que el proceso de eliminación de ruido realmente
reduce la entropía.

Antes de proceder con la demostración, realizaremos varias consideraciones sobre el
planteamiento del problema. En primer lugar, la fórmula de Tweedie \eqref{eq:tweedie} establece que
\begin{equation}
    \bar{\vx}^{\ast}(t, \vx_{t}) = \vx_{t} + t^{2}\nabla p_{t}(\vx_{t}),
\end{equation}
lo cual equipara un paso completo de eliminación de ruido desde el tiempo \(t\) hasta el tiempo \(0\) a un
paso de gradiente sobre la densidad logarítmica de \(\vx_{t}\). ¿Podemos obtener un resultado similar
para el paso completo de eliminación de ruido desde el tiempo \(t\) hasta el tiempo \(s\) en
\eqref{eq:app_denoising_iteration}? Resulta que sí podemos,
y es bastante sencillo. Utilizando \eqref{eq:app_denoising_iteration}
y la fórmula de Tweedie \eqref{eq:tweedie}, obtenemos
\begin{equation}\label{eq:app_denoising_iteration_score}
    \Ex[\vx_{s} \mid \vx_{t}] = \frac{s}{t}\vx_{t} + \bp{1 -
    \frac{s}{t}}\bp{\vx_{t} + t^{2}\nabla_{\vx_{t}}\log
    p_{t}(\vx_{t})} = \vx_{t} + \bp{1 -
    \frac{s}{t}}t^{2}\nabla_{\vx_{t}}\log p_{t}(\vx_{t}).
\end{equation}
Así pues, este paso iterativo de eliminación de ruido es, una vez más, un paso de gradiente sobre la
densidad logarítmica perturbada \(\log p_{t}\) con un tamaño de paso reducido. En
particular, este paso puede verse como una perturbación de la
distribución de la variable aleatoria \(\vx_{t}\) por el \textit{campo vectorial de la función de puntuación}, 
lo que sugiere una conexión con las ecuaciones diferenciales estocásticas
(SDE) y la teoría de los modelos de difusión
\cite{song2020score}. De hecho, se puede desarrollar una demostración del siguiente resultado
\Cref{thm:conditioning_reduces_entropy} utilizando este
potente marco y un argumento de límite (por ejemplo, siguiendo el
    enfoque técnico de la exposición de
\cite{DBLP:conf/iclr/ChenC0LSZ23}). Aquí daremos una prueba más sencilla,
que utilizará solo herramientas elementales y, por lo tanto, aclarará
algunas de las cantidades clave detrás del proceso de reducción de la entropía
mediante la eliminación de ruido. Por otro lado, tendremos que realizar algunos
cálculos ligeramente técnicos debido a que el proceso de eliminación de ruido
en \Cref{thm:conditioning_reduces_entropy} \textit{no}
se corresponde exactamente con el proceso \textit{inverso} asociado al
proceso de adición de ruido que genera la observación
\(\vx_{t}\).\footnote{Para quienes estén familiarizados con los modelos de difusión, nos
    referimos aquí al proceso directo con inversión temporal que no coincide con
    la secuencia de iteraciones generada por el proceso definido en el
    \Cref{thm:conditioning_reduces_entropy}. Estos procesos coinciden en
    un cierto límite en el que se dan infinitos pasos del
    \Cref{thm:conditioning_reduces_entropy} con niveles infinitamente
    pequeños de ruido añadidos en cada paso; para pasos generales y finitos,
    debemos introducir algunas aproximaciones independientemente del nivel de
sofisticación de nuestras herramientas.}

Queremos demostrar que \(h(\Ex[\vx_{s} \mid \vx_{t}]) < h(\vx_{t})\),
es decir, formalmente:
\begin{theorem}\label{thm:conditioning_reduces_entropy}
    Sea \(\vx\) una variable aleatoria cualquiera tal que los
    \Cref{assumption:entropy_x_compact_support,assumption:entropy_x_density}
    se cumplan, y sea \((\vx_{t})_{t \in [0, T]}\) el proceso estocástico
    \eqref{eq:app_additive_gaussian_noise_model}. Para cada
    \(t > 0\), escribamos
    \begin{equation}
        J(p_{t}) \doteq
        \int_{\R^{D}}\frac{\norm{\nabla
        p_{t}(\vxi)}_{2}^{2}}{p_{t}(\vxi)}\odif{\vxi}, \qquad
        U_{t} \doteq \max\rp{\frac{D}{t^{2}},
        \abs*{\frac{R^{2}}{t^{4}} - \frac{D}{t^{2}}}}
    \end{equation}
    para la información de Fisher de \(p_{t}\) y un límite uniforme
    sobre \(\abs{\Delta \log p_{t}}\), respectivamente. Entonces
    \begin{equation}
        h(\Ex[\vx_{s} \mid \vx_{t}]) < h(\vx_{t})
    \end{equation}
    para todo \(0 \leq s < t \leq T\) tal que
    \begin{equation}\label{eq:conditioning_entropy_condition}
        \bp{1 - \frac{s}{t}}t^{2}U_{t}^{2}
        \exp\bp{\bp{1 - \frac{s}{t}}t^{2}U_{t}} < 2J(p_{t}).
    \end{equation}
\end{theorem}
\begin{proof}
    Esta demostración se basa en dos ideas principales:
    \begin{enumerate}
        \item En primer lugar, calcula una densidad para \(\Ex[\vx_{s} \mid
            \vx_{t}]\) utilizando una fórmula de cambio de variables.
        \item En segundo lugar, limita esta densidad para controlar la entropía.
    \end{enumerate}
    El cambio de variables se justifica mediante \Cref{cor:gribonval_A2},
    que se derivó originalmente en \cite{Gribonval2011-pf}.

    A continuación, aplicamos estas ideas. A partir de \Cref{cor:gribonval_A2},
    obtenemos que la función \(\bar{\vx}\), definida como
    \(\bar{\vx}(\vxi) \doteq \Ex[\vx_{s} \given \vx_{t} = \vxi]\) es
    diferenciable, inyectiva y, por lo tanto, invertible en su imagen,
    que en adelante denotaremos como \(\cX \subseteq \R^{D}\). Denotamos
    su inversa como \(\bar{\vx}^{-1}\). Utilizando una fórmula de cambio de variables,
    la densidad \(\bar{p}\) de \(\bar{\vx}(\vx_{t})\) viene dada por
    \begin{equation}
        \bar{p}(\vxi) \doteq \frac{(p_{t} \circ
        \bar{\vx}^{-1})(\vxi)}{\det(\bar{\vx}^{\prime}(\bar{\vx}^{-1}(\vxi)))},
    \end{equation}
    donde (recordemos, según el \Cref{sec:autodiff}) \(\bar{\vx}^{\prime}\)
    es el jacobiano de \(\bar{\vx}\). Dado que por el
    \Cref{lem:gribonval_A1} sabemos que \(\bar{\vx}^{\prime}\) es una
    matriz definida positiva, el determinante es positivo y, por lo tanto, la
    densidad total es positiva. De ahí se deduce que
    \begin{align}
        h(\bar{\vx}(\vx_{t}))
        &= -\int_{\cX}\frac{(p_{t} \circ
        \bar{\vx}^{-1})(\vxi)}{\det(\bar{\vx}^{\prime}(\bar{\vx}^{-1}(\vxi)))}
        \log \frac{(p_{t} \circ
        \bar{\vx}^{-1})(\vxi)}{\det(\bar{\vx}^{\prime}(\bar{\vx}^{-1}(\vxi)))}
        \odif{\vxi} \\
        &= -\int_{\cX}\frac{(p_{t} \circ
        \bar{\vx}^{-1})(\vxi)}{\det(\bar{\vx}^{\prime}(\bar{\vx}^{-1}(\vxi)))}
        \log((p_{t} \circ \bar{\vx}^{-1})(\vxi))\odif{\vxi} \\
        &\qquad + \int_{\cX}\frac{(p_{t} \circ
        \bar{\vx}^{-1})(\vxi)}{\det(\bar{\vx}^{\prime}(\bar{\vx}^{-1}(\vxi)))}\log\det\rp{\bar{\vx}^{\prime}(\bar{\vx}^{-1}(\vxi))}\odif{\vxi}
        \\
        &= -\int_{\R^{D}}p_{t}(\vxi)\log p_{t}(\vxi)\odif{\vxi}
        + \int_{\cX}\frac{(p_{t} \circ
        \bar{\vx}^{-1})(\vxi)}{\det(\bar{\vx}^{\prime}(\bar{\vx}^{-1}(\vxi)))}\log\det\rp{\bar{\vx}^{\prime}(\bar{\vx}^{-1}(\vxi))}\odif{\vxi}
        \\
        &= h(\vx_{t}) - \int_{\cX}\frac{(p_{t} \circ
        \bar{\vx}^{-1})(\vxi)}{\det(\bar{\vx}^{\prime}(\bar{\vx}^{-1}(\vxi)))}\log\rp{\frac{1}{\det(\bar{\vx}^{\prime}(\bar{\vx}^{-1}(\vxi)))}}\odif{\vxi}.
    \end{align}
    Analizaremos el último término (incluido el \(-\)) y demostraremos que
    es negativo.

    Por la propiedad de la concavidad, se tiene que \(-x\log x \leq 1 - x\) para todo \(x \geq 0\). Por lo tanto
    \begin{align}
        h(\bar{\vx}(\vx_{t})) - h(\vx_{t})
        &= - \int_{\cX}\frac{(p_{t} \circ
        \bar{\vx}^{-1})(\vxi)}{\det(\bar{\vx}^{\prime}(\bar{\vx}^{-1}(\vxi)))}\log\rp{\frac{1}{\det(\bar{\vx}^{\prime}(\bar{\vx}^{-1}(\vxi)))}}\odif{\vxi}
        \\
        &\leq  \int_{\cX}(p_{t} \circ \bar{\vx}^{-1})(\vxi)\cdot
        \bp{1 -
        \frac{1}{\det(\bar{\vx}^{\prime}(\bar{\vx}^{-1}(\vxi)))}}\odif{\vxi} \\
        &= \int_{\cX}(p_{t} \circ \bar{\vx}^{-1})(\vxi)\odif{\vxi} -
        \int_{\cX}\frac{(p_{t} \circ
        \bar{\vx}^{-1})(\vxi)}{\det(\bar{\vx}^{\prime}(\bar{\vx}^{-1}(\vxi)))}\odif{\vxi}
        \\
        &=
        \int_{\R^{D}}p_{t}(\vxi)\det\rp{\bar{\vx}^{\prime}(\bar{\vx}^{-1}(\vxi))}\odif{\vxi}
        - \int_{\cX}\bar{p}(\vxi)\odif{\vxi} \\
        &= \int_{\R^{D}}p_{t}(\vxi)\det\rp{\vI + \bp{1 -
        \frac{s}{t}}t^{2}\nabla^{2}\log p_{t}(\vxi)}\odif{\vxi} - 1.
    \end{align}
    Ahora bien, por la desigualdad AM-GM para los valores propios, tenemos que, para cualquier
    matriz simétrica definida positiva \(\vM \in \PSD(D)\), la cota
    \begin{equation}
        \det(\vM)^{1/D} = \prod_{i = 1}^{D}\lambda_{i}(\vM)^{1/D}
        \leq \frac{\sum_{i = 1}^{D}\lambda_{i}(\vM)}{D} = \frac{\tr(\vM)}{D},
    \end{equation}
    que podemos aplicar a la expresión anterior y obtener
    \begin{align}
        &\int_{\R^{D}}p_{t}(\vxi)\det\rp{\vI + \bp{1 -
        \frac{s}{t}}t^{2}\nabla^{2}\log p_{t}(\vxi)}\odif{\vxi} \\
        &\leq \int_{\R^{D}} p_{t}(\vxi) \tr\rp{\frac{1}{D}\bs{\vI +
        \bp{1 - \frac{s}{t}}t^{2}\nabla^{2}\log p_{t}(\vxi)}}^{D}\odif{\vxi} \\
        &= \int_{\R^{D}} p_{t}(\vxi)\bp{1 + \frac{\bp{1 -
        \frac{s}{t}}t^{2}}{D}\tr(\nabla^{2}\log p_{t}(\vxi))}^{D}\odif{\vxi} \\
        &= \int_{\R^{D}} p_{t}(\vxi)\bp{1 + \frac{\bp{1 -
        \frac{s}{t}}t^{2}}{D}\Delta \log p_{t}(\vxi)}^{D}\odif{\vxi}.
    \end{align}
    Dado que \(\bar{\vx}^{\prime}(\vxi)\) es definida positiva
    (\Cref{cor:gribonval_A2}), tenemos que \(1 + \frac{(1 -
    s/t)t^{2}}{D}\Delta \log p_{t}(\vxi) =
    \frac{\tr(\bar{\vx}^{\prime}(\vxi))}{D} > 0\). Por la
    desigualdad \(\log(1 + x) \leq x\) para \(x > -1\),
    \begin{equation}
        \bp{1 + \frac{\bp{1 - \frac{s}{t}}t^{2}}{D}\Delta \log
        p_{t}(\vxi)}^{D} \leq \exp\bp{\bp{1 -
        \frac{s}{t}}t^{2}\Delta \log p_{t}(\vxi)}.
    \end{equation}
    Según el \Cref{lem:app_diffusion_laplacian_control} (donde, recordemos,
    \(R\) es el radio del soporte de \(\vx\), tal como se indica en el
    \Cref{assumption:entropy_x_compact_support}),
    \begin{equation}
        \abs{\Delta \log p_{t}(\vxi)} \leq \max\rp{\frac{D}{t^{2}},
        \abs*{\frac{R^{2}}{t^{4}} - \frac{D}{t^{2}}}} =: U_{t}.
    \end{equation}
    Si definimos \(\varepsilon \doteq \bp{1 -
    \frac{s}{t}}t^{2}U_{t}\), tenemos que \(\bp{1 -
    \frac{s}{t}}t^{2}\Delta \log p_{t}(\vxi) \leq \varepsilon\)
    para todo \(\vxi\). Por el teorema de Taylor con el resto de Lagrange,
    para cualquier \(c \leq \varepsilon\),
    \begin{equation}
        e^{c} \leq 1 + c + \frac{c^{2}}{2}e^{\varepsilon},
    \end{equation}
    puesto que \(e^{c} = 1 + c + \frac{c^{2}}{2}e^{\theta c}\) para algún
    \(\theta \in (0, 1)\), y
    \(e^{\theta c} \leq e^{\max(c, 0)} \leq e^{\varepsilon}\).
    Aplicando estos límites e integrando,
    \begin{align}
        &\int_{\R^{D}} p_{t}(\vxi)\bp{1 + \frac{\bp{1 -
        \frac{s}{t}}t^{2}}{D}\Delta \log
        p_{t}(\vxi)}^{D}\odif{\vxi} \\
        &\leq \int_{\R^{D}} p_{t}(\vxi)\exp\bp{\bp{1 -
        \frac{s}{t}}t^{2}\Delta \log
        p_{t}(\vxi)}\odif{\vxi} \\
        &\leq 1 + \bp{1 - \frac{s}{t}}t^{2}\int_{\R^{D}}
        p_{t}(\vxi)\Delta \log p_{t}(\vxi)\odif{\vxi}
        + \frac{e^{\varepsilon}}{2}\bp{1 -
        \frac{s}{t}}^{2}t^{4}\int_{\R^{D}}p_{t}(\vxi)\bp{\Delta
        \log p_{t}(\vxi)}^{2}\odif{\vxi} \\
        &= 1 - \bp{1 -
        \frac{s}{t}}t^{2}\int_{\R^{D}}\frac{\norm{\nabla
        p_{t}(\vxi)}_{2}^{2}}{p_{t}(\vxi)}\odif{\vxi}
        + \frac{e^{\varepsilon}}{2}\bp{1 -
        \frac{s}{t}}^{2}t^{4}\int_{\R^{D}}p_{t}(\vxi)\bp{\Delta
        \log p_{t}(\vxi)}^{2}\odif{\vxi},
    \end{align}
    donde la identidad \(\int p_{t}\Delta \log p_{t} = -\int
    \norm{\nabla p_{t}}_{2}^{2}/p_{t}\) es la misma que en la demostración
    del \Cref{thm:diffusion_entropy_increases}. Escribiendo \(J(p_{t})
    \doteq \int_{\R^{D}}\frac{\norm{\nabla
    p_{t}(\vxi)}_{2}^{2}}{p_{t}(\vxi)}\odif{\vxi}\) para la información de Fisher
    de \(p_{t}\), y limitando \(\bp{\Delta \log
    p_{t}}^{2} \leq U_{t}^{2}\), lo combinamos con nuestra estimación anterior
    para obtener
    \begin{equation}
        h(\bar{\vx}(\vx_{t})) - h(\vx_{t}) \leq -\bp{1 -
        \frac{s}{t}}t^{2}J(p_{t}) +
        \frac{e^{\varepsilon}\varepsilon^{2}}{2}.
    \end{equation}
    Esto es estrictamente negativo siempre que
    \(\varepsilon^{2}e^{\varepsilon} < 2\bp{1 -
    \frac{s}{t}}t^{2}J(p_{t})\). Sustituyendo
    \(\varepsilon = \frac{1 - \frac{s}{t}}t^{2}U_{t}\), esto es
    precisamente la condición
    \eqref{eq:conditioning_entropy_condition}.
\end{proof}

Nótese que la condición \eqref{eq:conditioning_entropy_condition}
siempre se cumple: para cualquier \(t > 0\) fijo, la información de Fisher
satisface \(J(p_{t}) > 0\) (puesto que \(p_{t}\) no es
constante) y \(U_{t} < \infty\), por lo que, si se elige un valor suficientemente pequeño de \(\bp{1 -
\frac{s}{t}}t^{2}\) lo suficientemente pequeño garantiza que
se cumpla \eqref{eq:conditioning_entropy_condition}.

Para comprender las implicaciones cuantitativas, consideremos el
comportamiento cuando \(t \to 0\) (cerca de los datos), donde la condición es
más restrictiva. Si se establece \(s = (1 - \eps)t\), la condición
se convierte en \(\eps t^{2}U_{t}^{2}\exp(\eps t^{2}U_{t}) <
2J(p_{t})\). Para valores pequeños de \(t\), el límite del laplaciano satisface
\(U_{t} \approx R^{2}/t^{4}\), por lo que el prefactor polinómico se escala como
\(\eps t^{2}U_{t}^{2} \approx \eps R^{4}/t^{6}\). El
escalado crítico es, por lo tanto,
\(\eps \sim t^{6}/R^{4}\). Si tomamos
\(\eps = \alpha t^{6}/R^{4}\),
la condición se reduce a \(\alpha < 2J(p_{t})\). En otras
palabras, cada paso de eliminación de ruido de tamaño
\(t - s = \eps\, t \sim t^{7}/R^{4}\) reduce la entropía.
Los pasos deben volverse progresivmente más finos a medida que \(t \to 0\), una
consecuencia del resto de segundo orden en la expansión de Taylor
utilizada en la demostración.

Al mismo tiempo, cabe señalar que es probable que nuestros límites sean imprecisos en
cuanto a la dependencia exacta de \(t\), y sin duda lo son para datos
con mayor estructura. Por ejemplo, el
crecimiento exponencial \(U_{t} \sim R^{2}/t^{4}\) refleja un límite del laplaciano en el peor de los casos
que puede ajustarse considerablemente bajo
supuestos adicionales de regularidad en \(p\). Además, si
\(p_{t}\) es log-cóncava (como ocurre, por ejemplo, cuando la propia \(p\) es
log-cóncava), entonces \(\Delta \log p_{t} \leq 0\) en todo el espacio, y se puede desarrollar un argumento considerablemente
más sencillo con conclusiones más sólidas.

\subsection{Lemas técnicos y resultados
auxiliares}\label{sub:app_diffusion_intermediate_results}

En esta subsección presentamos los resultados técnicos que fundamentan nuestros dos
teoremas conceptuales principales. Nuestra presentación seguirá más o menos
el enfoque habitual en matemáticas; empezaremos con los resultados de nivel más alto
y, progresivamente volveremos a los resultados más graduales. Los
resultados de mayor nivel utilizarán los resultados incrementales, y de esta
manera obtenemos un ordenamiento de dependencias legible y claro: ningún
resultado depende de los que le preceden. Los resultados que no dependen
unos de otros suelen ordenarse según el lugar en que aparecen en el
par de demostraciones anterior.

\subsubsection{Finitud de la entropía diferencial}

En primer lugar, demostramos que la entropía existe a lo largo del proceso estocástico
y que es finita.

\begin{lemma}\label{lem:diffusion_entropy_exists}
    Sea \(\vx\) una variable aleatoria cualquiera, y sea \((\vx_{t})_{t \in
    [0, T]}\) el proceso estocástico
    \eqref{eq:app_additive_gaussian_noise_model}.
    \begin{enumerate}
        \item Para \(t > 0\), la entropía diferencial \(h(\vx_{t})\)
            existe y es \(> -\infty\).
        \item Si, además, el
            \Cref{assumption:entropy_x_compact_support} se cumple para
            \(\vx\), entonces \(h(\vx) < \infty\) y \(h(\vx_{t})\ < \infty\).
    \end{enumerate}
\end{lemma}
\begin{proof}
    Para probar el \Cref{lem:diffusion_entropy_exists}.1, emplearemos un razonamiento analítico clásico
    y laborioso. Dado que \(\vx_{t}\) tiene una densidad, podemos escribir
    \begin{equation}
        h(\vx_{t}) = -\int_{\R^{D}}p_{t}(\vxi)\log p_{t}(\vxi) \odif{\vxi}.
    \end{equation}
    En consecuencia, definimos \(g \colon \R^{D} \to \R\) como
    \begin{equation}
        g(\vxi) \doteq -p_{t}(\vxi)\log p_{t}(\vxi) \implies
        h(\vx_{t}) = \int_{\R^{D}}g(\vxi)\odif{\vxi}.
    \end{equation}
    Como es habitual, para limitar el valor de una integral en análisis, se define
    \begin{eqnarray}
        &&g_{+}(\vxi) \doteq \max(g(\vxi), 0), \quad g_{-}(\vxi) \doteq
        \max(-g(\vxi), 0) \\ && \implies \quad g = g_{+} - g_{-}\quad
        \text{and} \quad g_{+}, g_{-} \geq 0.
    \end{eqnarray}
    Entonces
    \begin{equation}
        h(\vx_{t}) = \int_{\R^{D}}g_{+}(\vxi)\odif{\vxi} -
        \int_{\R^{D}}g_{-}(\vxi)\odif{\vxi},
    \end{equation}
    y se garantiza que ambas integrales son no negativas, ya que sus
    integrandos también lo son.

    Para demostrar que \(h(\vx_{t})\) está bien definida, necesitamos
    demostrar que \(\int_{\R^{D}}g_{+}(\vxi)\odif{\vxi} < \infty\) o
    \(\int_{\R^{D}}g_{-}(\vxi) \odif{\vxi} < \infty\). Para demostrar que
    \(h(\vx_{t}) > -\infty\), basta con demostrar que
    \(\int_{\R^{D}}g_{-}(\vxi)\odif{\vxi} < \infty\). Para acotar la
    integral de \(g_{-}\) necesitamos comprender la cantidad
    \(g_{-}\), es decir, queremos caracterizar cuándo \(g\) es negativo.
    \begin{equation}
        g(\vxi) \leq 0 \iff p_{t}(\vxi)\log p_{t}(\vxi) \geq 0 \iff
        \log p_{t}(\vxi) \geq 0 \iff p_{t}(\vxi) \geq 1.
    \end{equation}
    Por lo tanto, se cumple que
    \begin{equation}
        g_{-}(\vxi) = \indvar(p_{t}(\vxi) \geq 1)\cdot (-g(\vxi)) =
        \indvar(p_{t}(\vxi) \geq 1)p_{t}(\vxi)\log p_{t}(\vxi).
    \end{equation}
    Para poder acotar la integral de \(g_{-}(\vxi)\), necesitamos
    demostrar que \(p_{t}\) «no está demasiado concentrada», es decir, que
    \(p_{t}\) no es demasiado grande. Para demostrarlo, en este caso tenemos
    la suerte de poder acotar la propia función \(g_{-}(\vxi)\)
    en sí misma. Es decir, obsérvese que
    \begin{equation}
        \max_{\vxi \in \R^{D}}\phi_{t}(\vxi - \vx) = \phi_{t}(\vzero)
        = \frac{1}{(2\pi)^{D/2}t^{D}} =: C_{t}.
    \end{equation}
    que tiende a infinito cuando \(t \to 0\), pero es finito para todo \(t\) finito. Por lo tanto
    \begin{equation}
        p_{t}(\vxi) = \Ex \phi_{t}(\vxi - \vx) \leq \Ex C_{t} = C_{t}.
    \end{equation}
    Ahora hay dos casos.
    \begin{itemize}
        \item Si \(C_{t} < 1\), entonces \(p_{t}(\vxi) < 1\), por lo que el
            indicador nunca es \(1\); por lo tanto, \(g_{-} = 0\) en todos los casos
            y su integral también es \(0\).
        \item Si \(C_{t} \geq 1\), entonces \(\log C_{t} \geq 0\), por lo que
            dado que el logaritmo es monótonamente creciente,
            \begin{align}
                \int_{\R^{D}}g_{-}(\vxi)\odif{\vxi}
                &= \int_{\R^{D}}\indvar(p_{t}(\vxi) \geq
                1)p_{t}(\vxi)\log p_{t}(\vxi)\odif{\vxi}  \\
                &= \Ex[\indvar(p_{t}(\vx_{t}) \geq 1)\log p_{t}(\vx_{t})]  \\
                &\leq \Ex[\indvar(p_{t}(\vx_{t}) \geq 1) \log C_{t}] \\
                &= \Pr[p_{t}(\vx_{t}) \geq 1]\log C_{t}.
            \end{align}
    \end{itemize}
    Por lo tanto, tenemos que \(\int_{\R^{D}}g_{-}(\vxi)\odif{\vxi} < \infty\),
    por lo que la entropía diferencial \(h(\vx_{t})\) existe y es \(> -\infty\).

    Para demostrar el \Cref{lem:diffusion_entropy_exists}.2, supongamos que
    se cumple \Cref{assumption:entropy_x_compact_support}. Queremos
    demostrar que \(h(\vx) < \infty\) y \(h(\vx_{t}) < \infty\). El
    método para hacerlo es el mismo e implica el resultado de entropía máxima
    \Cref{thm:max_entropy}. Es decir, dado que \(\vx\) está
    absolutamente acotado, tiene una covarianza finita que
    denotaremos como \(\vSigma\). Entonces, la covarianza de \(\vx_{t}\) es
    \(\vSigma + t^{2}\vI\). Por lo tanto, la entropía de \(\vx\) y
    \(\vx_{t}\) puede tener como cota superior la entropía de las
    con las covarianzas respectivas, es decir, \(\log[(2\pi
    e)^{D}\det(\vSigma)]\) o \(\log[(2\pi e)^{D}\det(\vSigma +
    t^{2}\vI)]\), y ambas son \(< \infty\).
\end{proof}

\subsubsection{Integración por partes en la identidad de De Bruijn}

Por último, desarrollamos el argumento de la integración por partes al que se alude en
las demostraciones de los
\Cref{thm:diffusion_entropy_increases,thm:conditioning_reduces_entropy}.
El argumento es conceptualmente bastante simple, aunque exige algunas
estimaciones técnicas precisas para mostrar que el término de borde en la
integración por partes desaparece.

\begin{lemma}\label{lem:diffusion_ibp}
    Sea \(\vx\) una variable aleatoria tal que los
    \Cref{assumption:entropy_x_compact_support,assumption:entropy_x_density}
    se cumplen, y sea \((\vx_{t})_{t \in [0, T]}\) el proceso estocástico
    \eqref{eq:app_additive_gaussian_noise_model}. Para \(t
    \geq 0\), sea \(p_{t}\) la densidad de \(\vx_{t}\). Entonces, para una
    constante \(c \in \R\), se cumple
    \begin{equation}
        \int_{\R^{D}}\Delta p_{t}(\vxi)[c + \log
        p_{t}(\vxi)]\odif{\vxi} = -\int_{\R^{D}}\ip{\nabla \log
        p_{t}(\vxi)}{\nabla p_{t}(\vxi)}\odif{\vxi}.
    \end{equation}
\end{lemma}
\begin{proof}
    La idea básica de esta demostración se divide en dos pasos:
    \begin{itemize}
        \item En primer lugar, aplica el teorema de Green para realizar la integración por partes
            sobre un conjunto compacto;
        \item En segundo lugar, haz que el radio de este conjunto compacto tienda a
            \(+\infty\), para obtener integrales sobre todo el espacio \(\R^{D}\).
    \end{itemize}
    El teorema de Green establece que, para cualquier conjunto compacto \(\cK \subseteq
    \R^{D}\), una función \(\plainphi \colon
    \R^{D} \to \R\) dos veces continuamente diferenciable y una función \(\psi \colon
    \R^{D} \to \R\) continuamente diferenciable,
    \begin{equation}
        \int_{\cK}\bc{\psi(\vxi) \Delta \plainphi(\vxi) + \ip{\nabla
        \psi(\vxi)}{\nabla \plainphi(\vxi)}}\odif{\vxi} =
        \int_{\partial \cK}\psi(\vxi)\ip{\nabla
        \plainphi(\vxi)}{\vn(\vxi)}\odif{\sigma(\vxi)}
    \end{equation}
    donde \(\odif{\sigma(\vxi)}\) denota una integral sobre la
    «medida de superficie», es decir, la medida heredada en \(\partial
    \cK\), es decir, el contorno de \(\cK\), y, por consiguiente, \(\vxi\)
    toma valores en esta superficie y \(\vn(\vxi)\) es el vector normal unitario
    a \(\cK\) en el punto de la superficie \(\vxi\). Ahora, tomando
    \(\plainphi(\vxi) \doteq p_{t}(\vxi)\) y \(\psi(\vxi) \doteq c
    + \log p_{t}(\vxi)\), sobre una bola \(B_{r}(\vzero)\) de radio \ (r
    > 0) con centro en \(\vzero\) (de modo que \(\partial B_{r}(\vzero)\)
        sea la esfera de radio \(r\) con centro en \(\vzero\) y
    \(\vn(\vxi) = \vxi/\norm{\vxi}_{2} = \vxi/r\)):
    \begin{align}
        &\int_{B_{r}(\vzero)}\bc{\Delta p_{t}(\vxi)[c + \log
            p_{t}(\vxi)] + \ip{\nabla \log p_{t}(\vxi)}{\nabla
        p_{t}(\vxi)}}\odif{\vxi} \\
        &= \int_{\partial B_{r}(\vzero)}[c + \log
        p_{t}(\vxi)]\ip*{\nabla
        p_{t}(\vxi)}{\frac{\vxi}{r}}\odif{\sigma(\vxi)} \\
        &= \frac{1}{r}\int_{\partial B_{r}(\vzero)}[c + \log
        p_{t}(\vxi)]\ip*{\nabla p_{t}(\vxi)}{\vxi}\odif{\sigma(\vxi)}.
    \end{align}
    Al hacer que \(r \to \infty\), se cumple que
    \begin{align}
        &\int_{\R^{D}}\bc{\Delta p_{t}(\vxi)[c + \log p_{t}(\vxi)] +
        \ip{\nabla \log p_{t}(\vxi)}{\nabla p_{t}(\vxi)}}\odif{\vxi}
        \label{eq:app_diffusion_r_infinity_limit} \\
        &= \lim_{r \to \infty}\int_{B_{r}(\vzero)}\bc{\Delta
            p_{t}(\vxi)[c + \log p_{t}(\vxi)] + \ip{\nabla \log
        p_{t}(\vxi)}{\nabla p_{t}(\vxi)}}\odif{\vxi} \\
        &= \lim_{r \to \infty}\frac{1}{r}\int_{\partial
        B_{r}(\vzero)}[c + \log p_{t}(\vxi)]\ip*{\nabla
        p_{t}(\vxi)}{\vxi}\odif{\sigma(\vxi)},
    \end{align}
    donde la primera igualdad se deduce por convergencia dominada en
    el integrando. Queda por calcular el último límite. Para ello,
    tomamos expansiones asintóticas de cada término. El procedimiento principal es el
    siguiente: para \(\vxi \in \partial B_{r}(\vzero)\), tenemos
    \(\norm{\vxi}_{2} = r\), por lo que
    \begin{align}
        p_{t}(\vxi)
        &= \Ex[\phi_{t}(\vxi - \vx)] \\
        &= \Ex\rs{\underbrace{\frac{1}{(2\pi)^{D/2}t^{D}}}_{\doteq
        C_{t}}e^{-\|\vxi - \vx\|_{2}^{2}/(2t^{2})}} \\
        &= C_{t}\Ex\rs{e^{-(\norm{\vxi}_{2}^{2} - 2\ip{\vxi}{\vx} +
        \norm{\vx}_{2}^{2})/(2t^{2})}} \\
        &= C_{t}\Ex\rs{e^{-(r^{2} - 2\ip{\vxi}{\vx} +
        \norm{\vx}_{2}^{2})/(2t^{2})}} \\
        &= C_{t}e^{-r^{2}/(2t^{2})}\Ex[e^{(2\ip{\vxi}{\vx} -
        \norm{\vx}_{2}^{2})/(2t^{2})}].
    \end{align}
    Téngase en cuenta que, dado que \(\norm{\vxi}_{2} = r\), por el teorema de Cauchy-Schwarz se tiene que
    \begin{equation}
        -2r\norm{\vx}_{2} - \norm{\vx}_{2}^{2} \leq 2\ip{\vxi}{\vx} -
        \norm{\vx}_{2}^{2} \leq 2r\norm{\vx}_{2} - \norm{\vx}_{2}^{2}.
    \end{equation}
    Recordemos que, según el \Cref{assumption:entropy_x_compact_support},
    \(\vx\) tiene soporte en un conjunto compacto \(\cS\) de radio \(R\). Por lo tanto
    \begin{equation}
        -2R(r + R) \leq 2\ip{\vxi}{\vx} - \norm{\vx}_{2}^{2} \leq 2Rr.
    \end{equation}
    En otras palabras, se cumple que
    \begin{equation}
        C_{t}e^{-[r^{2} + 2R(r + R)]/(2t^{2})} \leq p_{t}(\vxi) \leq
        C_{t}e^{[-r^{2} + 2Rr]/(2t^{2})}.
    \end{equation}
    Ahora, para calcular el gradiente, podemos utilizar la
    \Cref{prop:p_t_derivatives} y la linealidad de la esperanza para calcular
    \begin{align}
        \ip{\nabla p_{t}(\vxi)}{\vxi}
        &= \ip*{-\frac{1}{t^{2}}\Ex\rs{\bp{\vxi - \vx}\phi_{t}(\vxi -
        \vx)}}{\vxi} \\
        &= -\frac{1}{t^{2}}\Ex\rs{\ip{\vxi - \vx}{\vxi}\phi_{t}(\vxi - \vx)} \\
        &= -\frac{1}{t^{2}}\Ex\rs{\bp{\norm{\vxi}_{2}^{2} -
        \ip{\vxi}{\vx}}\phi_{t}(\vxi - \vx)} \\
        &= -\frac{1}{t^{2}}\Ex\rs{\bp{r^{2} -
        \ip{\vxi}{\vx}}\phi_{t}(\vxi - \vx)} \\
        &= \frac{1}{t^{2}}\Ex\rs{\bp{\ip{\vxi}{\vx} -
        r^{2}}\phi_{t}(\vxi - \vx)}.
    \end{align}
    Utilizando de nuevo la desigualdad de Cauchy-Schwarz y la representación \(p_{t}(\vxi) \doteq
    \Ex[\phi_{t}(\vxi - \vx)]\), se cumple que
    \begin{align}
        &\frac{1}{t^{2}}\Ex\rs{\bp{-Rr - r^{2}}\phi_{t}(\vxi - \vx)}
        \leq \ip{\nabla p_{t}(\vxi)}{\vxi} \leq
        \frac{1}{t^{2}}\Ex\rs{\bp{Rr - r^{2}}\phi_{t}(\vxi - \vx)} \\
        &\frac{1}{t^{2}}\bp{-Rr - r^{2}}\Ex\rs{\phi_{t}(\vxi - \vx)}
        \leq \ip{\nabla p_{t}(\vxi)}{\vxi} \leq \frac{1}{t^{2}}\bp{Rr
        - r^{2}}\Ex\rs{\phi_{t}(\vxi - \vx)} \\
        &-\frac{r(R + r)}{t^{2}}p_{t}(\vxi) \leq \ip{\nabla
        p_{t}(\vxi)}{\vxi} \leq -\frac{r(r - R)}{t^{2}}p_{t}(\vxi).
    \end{align}
    Para \(r > R > 0\) (como corresponde, ya que vamos a tomar
    el límite \(r \to \infty\) mientras \(R\) es fijo), ambos lados son
    negativos. Esto tiene sentido: la mayor parte de la masa de probabilidad está
    contenida dentro de la bola de radio \(R\) y, por lo tanto, la puntuación
    apunta hacia adentro, teniendo un producto interno negativo con el
    vector que apunta hacia afuera \(\vxi\). Por lo tanto, utilizando las cotas apropiadas
    para \(p_{t}(\vxi)\),
    \begin{equation}
        -\frac{r(R + r)}{t^{2}}\cdot C_{t}e^{[-r^{2} + 2Rr]/(2t^{2})}
        \leq \ip{\nabla p_{t}(\vxi)}{\vxi} \leq -\frac{r(r -
        R)}{t^{2}}\cdot C_{t}e^{-[r^{2} + 2R(r + R)]/(2t^{2})}.
    \end{equation}
    Entonces, teniendo en cuenta que \(C_{t} = \mathrm{poly}(t^{-1})\), podemos calcular
    \begin{equation}
        [c + \log p_{t}(\vxi)]\ip{\nabla p_{t}(\vxi)}{\vxi} =
        \mathrm{poly}(r, R, t^{-1}, c)e^{-\Theta_{r}(r^{2})}
    \end{equation}
    Así pues, se puede ver que, si se establece que el área de la superficie de \(\partial
    B_{r}(\vzero)\) es \(\omega_{D} r^{D - 1}\), donde \(\omega_{D}\)
    es una función de \(D\), se cumple que
    \begin{equation}
        \frac{1}{r}\int_{\partial B_{r}(\vzero)}[c + \log
        p_{t}(\vxi)]\ip{\nabla p_{t}(\vxi)}{\vxi}\odif{\vxi} =
        \mathrm{poly}(r, R, t^{-1}, c)e^{-\Theta_{r}(r^{2})}
    \end{equation}
    y, por lo tanto, las colas que decaen exponencialmente significan
    \begin{equation}
        \lim_{r \to \infty}\frac{1}{r}\int_{\partial B_{r}(\vzero)}[c
        + \log p_{t}(\vxi)]\ip{\nabla p_{t}(\vxi)}{\vxi}\odif{\vxi} = 0.
    \end{equation}
    Por último, al sustituir en \eqref{eq:app_diffusion_r_infinity_limit}, obtenemos
    \begin{align}
        &\int_{\R^{D}}\bc{\Delta p_{t}(\vxi)[c + \log p_{t}(\vxi)] +
        \ip{\nabla \log p_{t}(\vxi)}{\nabla p_{t}(\vxi)}}\odif{\vxi} = 0 \\
        \implies
        &\int_{\R^{D}}\Delta p_{t}(\vxi)[c + \log
        p_{t}(\vxi)]\odif{\vxi} = -\int_{\R^{D}}\ip{\nabla \log
        p_{t}(\vxi)}{\nabla p_{t}(\vxi)}\odif{\vxi}
    \end{align}
    tal como se afirma.
\end{proof}

\subsubsection{Invertibilidad local del eliminador de ruido \(\bar{\vx}\)}

A continuación presentamos algunos resultados utilizados en la demostración del
\Cref{thm:conditioning_reduces_entropy}, que constituyen
generalizaciones adecuadas de los resultados correspondientes en \cite{Gribonval2011-pf}.

\begin{lemma}[Generalización de \cite{Gribonval2011-pf}, Lema
    A.1]\label{lem:gribonval_A1}
    Sea \(\vx\) una variable aleatoria tal que los
    \Cref{assumption:entropy_x_compact_support,assumption:entropy_x_density}
    se cumplen, y sea \((\vx_{t})_{t \in [0, T]}\) el proceso estocástico
    \eqref{eq:app_additive_gaussian_noise_model}. Sea \(s, t
    \in [0, T]\) tal que \(0 \leq s < t \leq T\), y sea
    \(\bar{\vx}(\vxi) \doteq \Ex[\vx_{s} \mid \vx_{t} = \vxi]\). El
    jacobiano \(\bar{\vx}^{\prime}(\vxi)\) es simétrico y definido positivo.
\end{lemma}
\begin{proof}
    Tenemos
    \begin{equation}
        \bar{\vx}^{\prime}(\vxi) = \vI + \bp{1 -
        \frac{s}{t}}t^{2}\nabla^{2}\log p_{t}(\vxi).
    \end{equation}
    Aquí expandimos
    \begin{equation}
        \nabla^{2}\log p_{t}(\vxi) =
        \frac{p_{t}(\vxi)\nabla^{2}p_{t}(\vxi) - (\nabla
        p_{t}(\vxi))(\nabla p_{t}(\vxi))^{\top}}{p_{t}(\vxi)^{2}}.
    \end{equation}
    Por lo tanto, debemos asegurarnos de que
    \begin{align}
        \bar{\vx}^{\prime}(\vxi)
        &= \vI + \bp{1 -
        \frac{s}{t}}t^{2}\frac{p_{t}(\vxi)\nabla^{2}p_{t}(\vxi) -
        (\nabla p_{t}(\vxi))(\nabla p_{t}(\vxi))^{\top}}{p_{t}(\vxi)^{2}} \\
        &= \frac{p_{t}(\vxi)^{2}\vI + \bp{1 -
            \frac{s}{t}}t^{2}\bs{p_{t}(\vxi)\nabla^{2}p_{t}(\vxi) -
        (\nabla p_{t}(\vxi))(\nabla p_{t}(\vxi))^{\top}}}{p_{t}(\vxi)^{2}}
    \end{align}
    es simétrica y semidefinida positiva. De hecho, es obviamente
    simétrica (por el teorema de Clairaut). Para demostrar su
    semidefinición positiva, sustituimos la representación de la esperanza de
    \(p_{t}\) dada por \eqref{eq:p_t_representation} (y \(\nabla
    p_{t}\), \(\Delta p_{t}\) por \Cref{prop:p_t_derivatives}) para
    obtener (donde \(\vx\) es como se definió y \(\vy\) es independiente e idénticamente distribuida como \(\vx\)),
    \begin{align}
        &\vv^{\top}[\bar{\vx}^{\prime}(\vxi)]\vv \\
        &= p_{t}(\vxi)^{-2}\vv^{\top}\Bigg\{p_{t}(\vxi)^{2}\vI +
            \bp{1 - \frac{s}{t}}t^{2}\Ex[\phi_{t}(\vxi -
            \vx)]\Ex\rs{\phi_{t}(\vxi - \vx)\cdot\frac{(\vxi -
            \vx)(\vxi - \vx)^{\top} - t^{2}\vI}{t^{4}}} \\
            & \qquad \qquad \qquad -\bp{1 -
            \frac{s}{t}}t^{2}\Ex\rs{-\phi_{t}(\vxi - \vx)\cdot\frac{\vxi
            - \vx}{t^{2}}}\Ex\rs{-\phi_{t}(\vxi - \vx)\cdot\frac{\vxi -
        \vx}{t^{2}}}^{\top}\Bigg\}\vv \\
        &= p_{t}(\vxi)^{-2}\vv^{\top}\Bigg\{\Ex[\phi_{t}(\vxi -
            \vx)\phi_{t}(\vxi - \vy)\vI] \\
            & \qquad \qquad \qquad + \bp{1 -
            \frac{s}{t}}t^{2}\Ex\rs{\phi_{t}(\vxi - \vx)\phi_{t}(\vxi
                - \vy)\cdot\frac{(\vxi - \vy)(\vxi - \vy)^{\top} -
            t^{2}\vI}{t^{4}}} \\
            & \qquad \qquad \qquad -\bp{1 -
            \frac{s}{t}}t^{2}\Ex\rs{\phi_{t}(\vxi - \vx)\phi_{t}(\vxi -
        \vy)\cdot\frac{(\vxi - \vx)(\vxi - \vy)^{\top}}{t^{4}}}\Bigg\} \\
        &= \frac{1 -
        s/t}{p_{t}(\vxi)^{2}}\vv^{\top}\Ex\mathopen{}\Bigg[\phi_{t}(\vxi
            - \vx)\phi_{t}(\vxi - \vy)\bc{\frac{1}{1 - s/t}\vI +
                \frac{(\vxi - \vy)(\vxi - \vy)^{\top}}{t^{2}} - \vI -
        \frac{(\vxi - \vx)(\vxi - \vy)^{\top}}{t^{2}}}\Bigg]\vv \\
        &= \frac{t - s}{tp_{t}(\vxi)^{2}}\vv^{\top}\Ex\rs{\frac{s}{t
            - s}\vI +\frac{(\vxi - \vx)(\vxi - \vx)^{\top}}{2t^{2}} +
            \frac{(\vxi - \vy)(\vxi - \vy)^{\top}}{2t^{2}} - \frac{(\vxi
        - \vx)(\vxi - \vy)^{\top}}{t^{2}}}\vv \\
        &= \frac{t - s}{tp_{t}(\vxi)^{2}}\vv^{\top}\Ex\rs{\frac{s}{t
            - s}\vI +\frac{1}{2t^{2}}\bp{(\vxi - \vx)(\vxi - \vx)^{\top}
                + (\vxi - \vy)(\vxi - \vy)^{\top} - 2(\vxi - \vx)(\vxi -
        \vy)^{\top}}}\vv \\
        &= \frac{t - s}{tp_{t}(\vxi)^{2}}\Ex\rs{\frac{s}{t -
            s}\norm{\vv}_{2}^{2} +\frac{1}{2t^{2}}\bp{[(\vxi -
                \vx)^{\top}\vv]^{2} + [(\vxi - \vy)^{\top}\vv]^{2} - 2[(\vxi
        - \vx)^{\top}\vv][(\vxi - \vy)^{\top}\vv]}} \\
        &= \frac{t - s}{tp_{t}(\vxi)^{2}}\Ex\rs{\frac{s}{t -
            s}\norm{\vv}_{2}^{2} +\frac{1}{2t^{2}}\bp{[(\vxi -
                \vx)^{\top}\vv]^{2} + [(\vxi - \vy)^{\top}\vv]^{2} - 2[(\vxi
        - \vx)^{\top}\vv][(\vxi - \vy)^{\top}\vv]}} \\
        &= \frac{t - s}{tp_{t}(\vxi)^{2}}\Ex\rs{\frac{s}{t -
            s}\norm{\vv}_{2}^{2} +\frac{1}{2t^{2}}\bp{[(\vxi -
        \vx)^{\top}\vv] - [(\vxi - \vy)^{\top}\vv]}^{2}} \\
        &= \frac{t - s}{tp_{t}(\vxi)^{2}}\Ex\rs{\frac{s}{t -
        s}\norm{\vv}_{2}^{2} +\frac{1}{2t^{2}}[(\vy - \vx)^{\top}\vv]^{2}} \\
        &= \frac{s}{tp_{t}(\vxi)^{2}}\norm{\vv}_{2}^{2} + \frac{t -
        s}{2t^{3}p_{t}(\vxi)}\Ex[\{(\vy - \vx)^{\top}\vv\}^{2}]
    \end{align}
    Dado que \(\vx\) y \(\vy\) son variables independientes e idénticamente distribuidas, la integral completa (es decir,
    la forma cuadrática original) es \(0\) si y solo si \(s = 0\)
    y \(\vx\) tiene un soporte totalmente contenido en un subespacio afín
    que es ortogonal a \(\vv\). Pero esto queda descartado por
    el supuesto (es decir, que \(\vx\) posee densidad en \(\R^{D}\)), por lo que
    el jacobiano \(\bar{\vx}^{\prime}(\vxi)\) es simétrico y definido positivo.
\end{proof}

\begin{lemma}[Generalización de \cite{Gribonval2011-pf} Corolario
    A.2, Parte 1]\label{lem:gribonval_A2}
    Sea \(f \colon \R^{D} \to \R^{D}\) una función diferenciable cualquiera
    cuyo jacobiano \(f^{\prime}(\vx)\) sea simétrico y definido positivo.
    Entonces \(f\) es inyectiva y, por lo tanto, invertible como
    función \(\R^{D} \to \Range(f)\), donde \(\Range(f)\) es el rango de \(f\).
\end{lemma}
\begin{proof}
    Supongamos que \(f\) no es inyectiva, es decir, que existen \(\vx,
    \vx^{\prime}\) tales que \(f(\vx) = f(\vx^{\prime})\) y \(\vx
    \neq \vx^{\prime}\). Definamos \(\vv \doteq (\vx^{\prime} -
    \vx)/\norm{\vx^{\prime} - \vx}_{2}\). Definamos la función \(g
    \colon \R \to \R\) como \(g(t) \doteq \vv^{\top}f(\vx + t\vv)\).
    Entonces \(g(0) = \vv^{\top}f(\vx) = \vv^{\top}f(\vx^{\prime}) =
    g(\norm{\vx^{\prime} - \vx}_{2})\). Dado que \(f\) es
    diferenciable, \(g\) es diferenciable, por lo que la derivada
    \(g^{\prime}\) debe ser nula para algún \(t^{\ast} \in (0,
    \norm{\vx^{\prime} - \vx}_{2})\) por el teorema del valor medio. Sin embargo,
    \begin{equation}
        g^{\prime}(t^{\ast}) \doteq \vv^{\top}\bs{f^{\prime}(\vx +
        t^{\ast}\vv)}\vv > 0
    \end{equation}
    ya que el jacobiano es definido positivo. Por lo tanto, llegamos a una
    contradicción, tal como se ha afirmado.
\end{proof}

A partir de la combinación de los dos resultados anteriores se obtiene el siguiente resultado fundamental.

\begin{corollary}[Generalización de \cite{Gribonval2011-pf} Corolario
    A.2, Parte 2]\label{cor:gribonval_A2}
    Sea \(\vx\) una variable aleatoria tal que los
    \Cref{assumption:entropy_x_compact_support,assumption:entropy_x_density}
    se cumplen, y sea \((\vx_{t})_{t \in [0, T]}\) el proceso estocástico
    \eqref{eq:app_additive_gaussian_noise_model}. Sea \(s, t
    \in [0, T]\) tales que \(0 \leq s < t \leq T\), y sea
    \(\bar{\vx}(\vxi) \doteq \Ex[\vx_{s} \mid \vx_{t} = \vxi]\). Entonces
    \(\bar{\vx}\) es inyectiva y, por lo tanto, invertible en su imagen.
\end{corollary}
\begin{proof}
    Lo único que queda por demostrar es que \(\bar{\vx}\) es
    diferenciable, pero esto se deduce inmediatamente de la fórmula de Tweedie
    (\Cref{thm:tweedie}), que muestra que \(\bar{\vx}\) es
    diferenciable si y solo si \(\nabla \log p_{t}\) es
    diferenciable, y esto viene dado por la \Cref{eq:p_t_representation}.
\end{proof}

\subsubsection{Control del laplaciano \(\Delta \log p_{t}\)}

Finalmente, desarrollamos una estimación técnica requerida para 
probar el \Cref{thm:conditioning_reduces_entropy} lo cual motiva directamente
el supuesto sobre los valores admisibles de \(t\).

\begin{lemma}\label{lem:app_diffusion_laplacian_control}
    Sea \(\vx\) una variable aleatoria tal que los
    \Cref{assumption:entropy_x_compact_support,assumption:entropy_x_density}
    se cumplan, y sea \((\vx_{t})_{t \in [0, T]}\) el proceso estocástico
    \eqref{eq:app_additive_gaussian_noise_model}. Sea
    \(p_{t}\) la densidad de \(\vx_{t}\). Entonces, para \(t > 0\) se cumple
    \begin{equation}
        \sup_{\vxi \in \R^{D}}\abs{\Delta \log p_{t}(\vxi)} \leq
        \max\rp{\frac{D}{t^{2}}, \abs*{\frac{R^{2}}{t^{4}} - \frac{D}{t^{2}}}}.
    \end{equation}
\end{lemma}
\begin{proof}
    Mediante la regla de la cadena, un ejercicio sencillo permite calcular
    \begin{equation}
        \Delta \log p_{t}(\vxi) = \frac{\Delta
        p_{t}(\vxi)}{p_{t}(\vxi)} - \frac{\norm{\nabla
        p_{t}(\vxi)}_{2}^{2}}{p_{t}(\vxi)^{2}}.
    \end{equation}
    Utilizando la \Cref{prop:p_t_derivatives} para expresar los términos de \(\Delta
    p_{t}(\vxi)\), obtenemos
    \begin{align}
        \frac{\Delta p_{t}(\vxi)}{p_{t}(\vxi)}
        &= \frac{\Ex\rs{\frac{\norm{\vxi - \vx}_{2}^{2} -
        Dt^{2}}{t^{4}} \cdot \phi_{t}(\vxi - \vx)}}{\Ex[\phi_{t}(\vxi
        - \vx)]} \\
        &= \frac{\int_{\R^{D}}\bc{\frac{\norm{\vxi - \vu}_{2}^{2} -
            Dt^{2}}{t^{4}}}\phi_{t}(\vxi -
        \vu)p(\vu)\odif{\vu}}{\int_{\R^{D}}\phi_{t}(\vxi -
        \vu)p(\vu)\odif{\vu}}.
    \end{align}
    Esto parece una marginalización bayesiana, así que definamos la
    densidad normalizada adecuada
    \begin{equation}
        q_{\vxi}(\vu) = \frac{\phi_{t}(\vxi -
        \vu)p(\vu)}{\int_{\R^{D}}\phi_{t}(\vxi -
        \vv)p(\vv)\odif{\vv}} = \frac{\phi_{t}(\vxi -
        \vu)p(\vu)}{\Ex[\phi_{t}(\vxi - \vx)]} = \frac{\phi_{t}(\vxi
        - \vu)p(\vu)}{p_{t}(\vxi)}
    \end{equation}
    Entonces, definiendo \(\vy_{\vxi} \sim q_{\vxi}\), podemos escribir
    \begin{equation}
        \frac{\Delta p_{t}(\vxi)}{p_{t}(\vxi)} =
        \int_{\R^{D}}\bc{\frac{\norm{\vxi - \vu}_{2}^{2} -
        Dt^{2}}{t^{4}}}q_{\vxi}(\vu)\odif{\vu} =
        \frac{1}{t^{4}}\Ex[\norm{\vxi - \vy_{\vxi}}_{2}^{2}] - \frac{D}{t^{2}}.
    \end{equation}
    Del mismo modo, al desarrollar el segundo término (sin elevar al cuadrado) obtenemos
    \begin{equation}
        \frac{\nabla p_{t}(\vxi)}{p_{t}(\vxi)} = -\frac{\vxi -
        \Ex[\vy_{\vxi}]}{t^{2}}.
    \end{equation}
    Si se establece que \(\vz_{\vxi} \doteq \vy_{\vxi} - \vxi\), se cumple que
    \begin{equation}
        \frac{\Delta p_{t}(\vxi)}{p_{t}(\vxi)} =
        \frac{\Ex[\norm{\vz_{\vxi}}_{2}^{2}]}{t^{4}} -
        \frac{D}{t^{2}}, \qquad \frac{\nabla
        p_{t}(\vxi)}{p_{t}(\vxi)} = \frac{\Ex[\vz_{\vxi}]}{t^{2}}.
    \end{equation}
    Por lo tanto, al desarrollar completamente \(\Delta \log p_{t}\), tenemos
    \begin{align}
        \Delta \log p_{t}(\vxi)
        &= \frac{\Ex[\norm{\vz_{\vxi}}_{2}^{2}]}{t^{4}} -
        \frac{D}{t^{2}} - \frac{\norm{\Ex[\vz_{\vxi}]}_{2}^{2}}{t^{4}} \\
        &= \frac{\Ex[\norm{\vz_{\vxi}}_{2}^{2}] -
        \norm{\Ex[\vz_{\vxi}]}_{2}^{2}}{t^{4}} - \frac{D}{t^{2}} \\
        &= \frac{\tr(\Cov(\vz_{\vxi}))}{t^{4}} - \frac{D}{t^{2}} \\
        &= \frac{\tr(\Cov(\vy_{\vxi}))}{t^{4}} - \frac{D}{t^{2}}.
    \end{align}
    Una cota inferior trivial de esta traza es \(0\), ya que las matrices de covarianza
    son semidefinidas positivas. Para hallar una cota superior, obsérvese
    que \(\vy_{\vxi}\) toma valores únicamente en el soporte de \(\vx\)
    (puesto que \(p\) es un factor de la densidad \(q_{\vxi}\) de
    \(\vy_{\vxi}\)), el cual, por el
    \Cref{assumption:entropy_x_compact_support}, es un conjunto compacto
    \(\cS\) con radio \(R \doteq \sup_{\vxi \in \cS}\norm{\vxi}_{2}\). Por lo tanto
    \begin{equation}
        \tr(\Cov(\vy_{\vxi})) = \Ex[\norm{\vy_{\vxi}}_{2}^{2}] -
        \norm{\Ex[\vy_{\vxi}]}_{2}^{2} \leq
        \Ex[\norm{\vy_{\vxi}}_{2}^{2}] \leq R^{2}.
    \end{equation}
    Por lo tanto
    \begin{equation}
        -\frac{D}{t^{2}} \leq \Delta \log p_{t}(\vxi) \leq
        \frac{R^{2}}{t^{4}} - \frac{D}{t^{2}},
    \end{equation}
    que muestra lo que se sostiene.
\end{proof}

\subsubsection{Cálculos de derivadas}

A continuación presentamos algunas derivaciones útiles que serán 
reutilizadas en todo el apéndice.

\begin{proposition}\label{prop:p_t_derivatives}
    Sea \(\vx\) una variable aleatoria cualquiera tal que los
    \Cref{assumption:entropy_x_compact_support,assumption:entropy_x_density}
    se cumplan, y sea \((\vx_{t})_{t \in [0, T]}\) el proceso estocástico
    \eqref{eq:app_additive_gaussian_noise_model}. Para \(t
    \geq 0\), sea \(p_{t}\) la densidad de \(\vx_{t}\). Entonces
    \begin{align}
        \pdv{p_{t}}{t}(\vxi)
        &= \Ex\rs{\phi_{t}(\vxi - \vx)\cdot\frac{\norm{\vxi -
        \vx}_{2}^{2} - Dt^{2}}{t^{3}}} \\
        \nabla p_{t}(\vxi)
        &= -\Ex\rs{\phi_{t}(\vxi - \vx)\cdot\frac{\vxi - \vx}{t^{2}}} \\
        \nabla^{2}p_{t}(\vxi)
        &= \Ex\rs{\phi_{t}(\vxi - \vx)\cdot\frac{(\vxi - \vx)(\vxi -
        \vx)^{\top} - t^{2}\vI}{t^{4}}} \\
        \Delta p_{t}(\vxi)
        &= \Ex\rs{\phi_{t}(\vxi - \vx)\cdot\frac{\norm{\vxi -
        \vx}_{2}^{2} - Dt^{2}}{t^{4}}}.
    \end{align}
\end{proposition}
\begin{proof}
    Utilizamos la representación por convolución de \(p_{t}\), es decir,
    \eqref{eq:p_t_representation}. Tomando primero la derivada temporal,
    un cálculo muestra que se aplica la \Cref{prop:dutis},\footnote{Utilizamos
        \(f_{t}(\vxi) = p(\vxi) \phi_{t}(\vxi - \vx)\), observando que
        es dos veces continuamente diferenciable en \(\vxi\) y (más
        que dos veces en \(t\). A continuación, para verificar
        la integrabilidad local de \(f_{t}\), calculamos
        \(\pdv{f_{t}}{t}(\vxi) = f_{t}(\vxi)\cdot
        \frac{1}{t^{3}}(\norm{\vxi - \vx}_{2}^{2} - Dt^{2})\), lo cual
        es fácil de comprobar que es integrable sobre \(\vxi\) y \(t \in [t_{\min},
        t_{\max}]\) donde \(t_{\min} > 0\). (De hecho, \(f_{t}\) tiene
            colas que decaen exponencialmente, por lo que el término cuadrático en el
    producto no supone ningún problema.)} así que podemos llevar la derivada al interior
    de la integral/esperanza como:
    \begin{equation}
        \pdv{p_{t}}{t}(\vxi) = \pdv*{\Ex[\phi_{t}(\vxi - \vx)]}{t} =
        \Ex\rs{\pdv*{\phi_{t}(\vxi - \vx)}{t}} = \pdv{\phi_{t}}{t} * p.
    \end{equation}
    Por otra parte, gracias a las propiedades de las convoluciones
    (\Cref{prop:diff_convolution}) y utilizando el hecho de que \(p\) tiene
    soporte compacto (\Cref{assumption:entropy_x_compact_support}),
    \begin{equation}
        p_{t} = \phi_{t} * p \implies \nabla p_{t} = \nabla \phi_{t}
        * p \implies \nabla^{2}p_{t} = \nabla^{2}\phi_{t} * p
        \implies \Delta p_{t} = \Delta \phi_{t} * p.
    \end{equation}
    El resto del cálculo se deduce de la \Cref{prop:normal_derivatives}.
\end{proof}

\begin{proposition}\label{prop:normal_derivatives}
    Para \(t > 0\) y \(\vxi \in \R^{D}\), se cumple que
    \begin{align}
        \pdv*{\phi_{t}}{t}(\vxi)
        &= \phi_{t}(\vxi) \cdot \frac{\norm{\vxi}_{2}^{2} - Dt^{2}}{t^{3}} \\
        \nabla \phi_{t}(\vxi)
        &= -\phi_{t}(\vxi)\cdot\frac{\vxi}{t^{2}} \\
        \nabla^{2} \phi_{t}(\vxi)
        &= \phi_{t}(\vxi)\cdot\frac{\vxi\vxi^{\top} - t^{2}\vI}{t^{4}} \\
        \Delta \phi_{t}(\vxi)
        &= \phi_{t}(\vxi) \cdot \frac{\norm{\vxi}_{2}^{2} - Dt^{2}}{t^{4}}.
    \end{align}
\end{proposition}
\begin{proof}
    Cálculo directo.
\end{proof}

\subsubsection{Derivación bajo el signo integral}

En este apéndice, derivamos muchas
veces bajo el signo de la integral, y resulta esencial conocer las condiciones 
bajo las cuales es posible proceder. Hay dos
tipos de derivación bajo el signo de la integral:
\begin{enumerate}
    \item Derivación de la integral \(\int f_{t}(\vxi)\odif{\vxi}\)
        con respecto al parámetro auxiliar \(t\).
    \item Derivación de una convolución \((f * g)(\vxi) = \int
        f(\vxi)g(\vxi - \vu)\odif{\vu}\) con respecto a la variable \(\vxi\).
\end{enumerate}

Para la primera categoría, presentamos un resultado explícito, enunciado sin 
prueba pero atribuible a 
\href{https://planetmath.org/differentiationundertheintegralsign}{la
referencia vinculada}, que deduce el siguiente resultado como un caso especial
de un teorema más general sobre la interacción de
operadores diferenciales y distribuciones templadas, lo cual excede con creces el alcance del
libro. Se puede encontrar una referencia formal completa en \cite{jones1982theory}.
\begin{proposition}[\cite{jones1982theory}, Sección 11.12]\label{prop:dutis}
    Sea \(f \colon (0, T) \times \R^{D} \to \R\) tal que:
    \begin{itemize}
        \item \(f\) es una función conjuntamente medible de \((t, \vxi)\);
        \item Para casi todo \(\vxi \in \R^{D}\) en el sentido de Lebesgue, la
            función \(t \mapsto f_{t}(\vxi)\) es absolutamente continua;
        \item \(\pdv{f_{t}}{t}\) es localmente integrable, es decir, para
            cada \([t_{\min}, t_{\max}] \subseteq (0, T)\) se cumple que
            \begin{equation}
                \int_{t_{\min}}^{t_{\max}}\int_{\R^{D}}\abs*{\pdv{f_{t}}{t}(\vxi)}\odif{\vxi}
                < \infty.
            \end{equation}
    \end{itemize}
    Entonces, \(t \mapsto \int_{\R^{D}}f_{t}(\vxi)\odif{\vxi}\) es una
    función absolutamente continua en \((0, T)\), y su derivada es
    \begin{equation}
        \odv*{\int_{\R^{D}}f_{t}(\vxi)\odif{\vxi}}{t} =
        \int_{\R^{D}}\pdv*{f_{t}}{t}(\vxi)\odif{\vxi},
    \end{equation}
    definida para casi todo \(t \in (0, T)\).
\end{proposition}

Para la segunda categoría, ofrecemos otro resultado concreto, que se expone
sin demostración, pero que está completamente formalizado en \cite{brezis2011functional}.
\begin{proposition}[\cite{brezis2011functional}, Proposición
    4.20]\label{prop:diff_convolution}
    Sea \(f\) una función \(k\)-veces continuamente diferenciable con soporte compacto,
    y sea \(g\) una función localmente integrable. Entonces, la
    convolución \(f * g\) definida por
    \begin{equation}
        (f * g)(\vxi) \doteq \int_{\R^{D}}f(\vu)g(\vxi - \vu)\odif{\vu}
    \end{equation}
    es \(k\)-veces continuamente diferenciable, y su derivada de
    orden \(k\) es
    \begin{equation}
        \nabla^{k}(f * g) =(\nabla^{k}f) * g.
    \end{equation}
\end{proposition}
Aunque no aparece en el libro, un sencillo razonamiento basado en la integración por partes
demuestra que, si \(g\) también es \(k\)-veces diferenciable, entonces podemos
«intercambiar» la regularidad:
\begin{equation}
    \nabla^{k}(f * g) = f * (\nabla^{k} g).
\end{equation}

\section{Codificación con pérdida y empaquetamiento esférico}\label{app:rate-distortion-covering}

En esta sección, demostramos el \Cref{thm:covering-number-rate-distortion}.
Siguiendo las convenciones utilizadas a lo largo de este apéndice, escribimos $\cS = \Supp(\vx)$
para el soporte compacto de la variable aleatoria $\vx$.

Tal como se ha adelantado, haremos un supuesto de regularidad sobre el conjunto de soporte $\cS$
para demostrar el \Cref{thm:covering-number-rate-distortion}. Una posibilidad para
proceder bajo supuestos mínimos sería aplicar los resultados de
\cite{Riegler2018-jh,Riegler2023-rr} en nuestro contexto, ya que estos resultados se aplican
a conjuntos $\cS$ con muy baja regularidad (por ejemplo, conjuntos tipo Cantor con estructura
fractal). Sin embargo, hemos encontrado que el cálculo preciso de estas constantes 
---una tarea necesaria para demostrar una conclusión como el
\Cref{thm:covering-number-rate-distortion}---resulta bastante tedioso en nuestro
contexto.
Por lo tanto, nuestro enfoque consiste en añadir un supuesto de regularidad geométrica sobre
el conjunto $\cS$ que cede en términos de generalidad, pero nos permite desarrollar
un argumento más transparente. Para evitar sacrificar demasiada generalidad, debemos
asegurarnos de que no se prohíba la baja dimensionalidad en el conjunto $\cS$.
Por lo tanto, consideramos el ejemplo que hemos utilizado a lo largo del libro: la
mezcla de distribuciones gaussianas de rango bajo. En este marco geométrico,
modelamos $\cS$ como una unión de hiperesferas, cada una de dimensión $d_k$ (posiblemente
mucho menor que $D$), situadas en subespacios mutuamente ortogonales de $\R^D$.
Esta es una simplificación geométrica del modelo de mezcla gaussiana: el soporte de cada
componente se aproxima mediante una esfera en el subespacio generado por sus
direcciones principales. Cuando las dimensiones de las componentes $d_k$ son grandes, este
supuesto equivale a un supuesto de mezcla de gaussianas de bajo rango, debido a
la concentración de medidas en espacios de alta dimensión. Los modelos CRATE que desarrollamos y entrenamos en los
\Cref{ch:deep-networks,ch:applications} tienen sus \textit{dimensiones de subespacio}
establecidas de manera compatible con este supuesto.

\begin{assumption}\label{assumption:union-of-spheres}
    El soporte $\cS \subset \R^D$ de la variable aleatoria $\vx$ es una unión finita
    de $K$ esferas, cada una de dimensión $d_k$, donde $k \in [K]$.
    La probabilidad de que $\vx$ se extraiga de la $k$-ésima esfera viene dada por
    $\pi_k \in [0, 1]$, y, condicionalmente a que se extraiga de la $k$-ésima esfera,
    $\vx$ se distribuye uniformemente en esa esfera.
    Los soportes satisfacen que cada esfera es mutuamente ortogonal con todas
    las demás.
\end{assumption}

Procedemos bajo el \Cref{assumption:union-of-spheres} simplificadora con el fin de
evitar un exceso de tecnicismo y establecer una conexión con un importante ejemplo
que se utiliza a lo largo de toda la monografía. Creemos que nuestros resultados pueden generalizarse
para abarcar $\cS$ de la clase de \textit{conjuntos con alcance positivo} con
un esfuerzo técnico adicional, pero dejamos esto para el futuro.

\subsection{Demostración de la relación entre la distorsión de la tasa y la cobertura}

Esbozamos brevemente la demostración, luego pasamos a establecer tres
lemas fundamentales y, a continuación, presentamos la demostración. La demostración se basará en los conceptos introducidos en el
esbozo que figura a continuación.

Obtener un límite superior para la función de tasa-distorsión
\eqref{eqn:rate-distortion-general} es sencillo: según la caracterización de la tasa
(es decir, la función de tasa-distorsión es la tasa mínima de
un código para $\vx$ con una distorsión cuadrática esperada $\ell^2$ de $\epsilon$), para establecer un límite superior
de $\cR_{\epsilon}(\x)$ solo es necesario demostrar que existe un código para $\ x$ que
alcance esta distorsión objetivo, y cualquier cobertura $\epsilon$ de $\Supp(\x)$
lo logra, con una tasa igual al logaritmo en base 2 de la cardinalidad de
la cobertura.
El límite inferior es más sutil. Utilizamos el límite inferior de Shannon,
discutido en \Cref{rem:slb}: al calcular las constantes en \cite[\S III,
(22)]{Linder1994-ej} se obtiene una versión más precisa del resultado citado en la
\Cref{eq:slb} (en bits, por supuesto): para cualquier variable aleatoria $\vx$ con soporte compacto
y una densidad, se cumple que
\begin{equation}\label{eq:slb-appendix-1}
    \cR_{\epsilon}(\x)
    \geq
    h(\x)
    - \log \volume(B_{\epsilon})
    +
    \log
    \left(
        \frac{
            2
        }
        {
            D \Gamma(D/2)
        }
        \left(
            \frac{
                D
            }
            {
                2e
            }
        \right)^{D/2}
    \right),
\end{equation}
con la entropía (etc.) en nats en esta expresión.
La constante se puede estimar fácilmente utilizando la aproximación de Stirling.
Una forma cuantitativa de la aproximación de Stirling que suele resultar útil da
para cualquier $x > 0$ \cite{Jameson2015-hy}
\begin{equation}
    \Gamma(x) \leq \sqrt{2\pi} x^{x - 1/2} e^{-x} e^{1/(12x)}.
\end{equation}
Aplicaremos esta cota a $\Gamma(D/2)$ en la \Cref{eq:slb-appendix-1}.
Se obtiene
\begin{align}
    \log
    \left(
        \frac{
            2
        }
        {
            D \Gamma(D/2)
        }
        \left(
            \frac{
                D
            }
            {
                2e
            }
        \right)^{D/2}
    \right)
    &\geq
    -\frac{1}{6D}
    +
    \log\left(
        \frac{2}{D}
        \left(
            \frac{
                D
            }
            {
                2e
            }
        \right)^{D/2}
        \cdot
        \sqrt{\frac{D}{4\pi}}
        \left(
            \frac{D}{2e}
        \right)^{-D/2}
    \right)
    \\
    &=
    -\frac{1}{6D}
    - \frac{1}{2}\log D\pi,\label{eq:slb-constant-est}
\end{align}
que podemos tomar como el valor explícito de la constante $C_D$ en \Cref{eq:slb}.
En resumen, el límite inferior de Shannon totalmente cuantificado (en bits):
\begin{equation}\label{eq:slb-appendix}
    \cR_{\epsilon}(\x)
    \geq
    h(\x)
    - \log_2 \volume(B_{\epsilon})
    - O(\log D).
\end{equation}

Ahora bien, la restricción importante para nuestros propósitos actuales es que el límite inferior de Shannon
exige que la variable aleatoria $\vx$ tenga una densidad, lo que descarta muchas
distribuciones de interés de baja dimensión.
Pero consideremos por un momento la situación en la que $\vx$ sí admite una densidad.
El supuesto de que $\vx$ sigue una distribución uniforme en su soporte se formaliza fácilmente
en este contexto: para cualquier conjunto de Borel $A \subset \cS$, tenemos
\begin{equation}
    \Pr[\vx \in A] = \int_{A} \frac{1}{\volume(\cS)} \odif \vx.
\end{equation}
Entonces, la entropía $h(\vx)$ es simplemente
\begin{equation}
    h(\vx) = \log_2 \volume(\cS).
\end{equation}
La demostración concluye entonces con un lema que relaciona la razón $\volume(\cS)
/ \volume(B_{\epsilon})$ con el número de recubrimiento $\epsilon$ de $\cS$ mediante
bolas de tamaño $\epsilon$.

Para ampliar el programa anterior a distribuciones degeneradas que satisfagan el
\Cref{assumption:union-of-spheres},
nuestra demostración del límite inferior en \Cref{thm:covering-number-rate-distortion}
aprovechará un argumento de aproximación de la distribución real de baja dimensión
$\vx$ mediante distribuciones «cercanas» que tienen
densidades, de manera similar pero no exactamente igual a la esbozada en la demostración
que precede al \Cref{thm:max_entropy}.
A continuación, relacionaremos el parámetro introducido en
la secuencia de aproximación con el parámetro de distorsión $\epsilon$ con el fin de
obtener la conclusión deseada en el \Cref{thm:covering-number-rate-distortion}.

\begin{definition}\label{def:thickening-set}
    Sea $\cS$ un conjunto compacto.
    Para cualquier $\delta > 0$, definimos el entorno tubular de $\cS$ de espesor $\delta$, denotado
    $\cS_{\delta}$, como
    \begin{equation}
        \cS_{\delta} = \set*{\vxi \in \R^D \given \mathrm{dist}(\vxi, \cS) \leq
        \delta}.
    \end{equation}
\end{definition}
La función de distancia a la que se hace referencia en la \Cref{def:thickening-set} se define mediante
\begin{equation}\label{eq:dist-func}
    \dist(\vxi, \cS) = \inf_{\vxi' \in \cS} \norm*{\vxi - \vxi'}_2.
\end{equation}
Para un conjunto compacto $\cS$, el teorema de Weierstrass implica que, para cualquier $\vxi \in
\R^D$, siempre existe algún $\vxi' \in \cS$ que alcanza el infimo en
la función de distancia.
La compacidad de $\cS_{\delta}$ se deduce fácilmente de la compacidad de $\cS$, por lo que
$\volume(\cS_{\delta})$ es finito para cualquier $\delta > 0$. Entonces es posible
establecer la siguiente definición de una variable aleatoria engrosada, especializada al
\Cref{assumption:union-of-spheres}.

\begin{definition}\label{def:thickening-rv-uos}
    Sea $\vx$ una variable aleatoria tal que $\Supp(\vx) = \cS$ es una unión de
    $K$ hiperesferas, distribuidas como en el \Cref{assumption:union-of-spheres}.
    Denotemos por $\cS_k$ el soporte de cada componente de la mezcla.
    Definamos la variable aleatoria engrosada $\vx_{\delta}$ como la mezcla de
    medidas en la que cada medida componente es uniforme en el conjunto engrosado
    $\cS_{k, \delta}$ (\Cref{def:thickening-set}), para $k \in [K]$, con pesos de mezcla
    $\pi_k$.
\end{definition}

\begin{lemma}\label{lem:rate-distortion-lb-uos}
    Supongamos que la variable aleatoria $\vx$ satisface el
    \Cref{assumption:union-of-spheres}. Entonces, si $0 < \delta < \tfrac{1}{2}$, la
    variable aleatoria engrosada $\vx_{\delta}$ (\Cref{def:thickening-rv-uos})
    satisface, para cualquier $\epsilon > 0$
    \begin{equation}
        R_{\delta + \epsilon}(\vx_{\delta})
        \leq
        R_{\epsilon}(\vx).
    \end{equation}
\end{lemma}

La demostración del \Cref{lem:rate-distortion-lb-uos} se pospone hasta
la \Cref{sec:app-rate-dist-deferred-proofs}.
Utilizando el \Cref{lem:rate-distortion-lb-uos}, el programa anterior puede llevarse a cabo,
ya que la variable aleatoria $\vx_{\delta}$ tiene una densidad uniforme
con respecto a la medida de Lebesgue.

\begin{proof}[(Demostración del \Cref{thm:covering-number-rate-distortion})]
    La cota superior se demuestra fácilmente. Si $S$ es cualquier recubrimiento $\epsilon$ del
    soporte de $\vx$ con cardinalidad $\cN_{\epsilon}(\Supp(\vx))$,
    entonces consideremos el
    esquema de codificación que asigna a cada $\vxi \in \Supp(\vx)$ la reconstrucción
    $\hat{\vxi} = \argmin_{\vxi' \in S}\, \norm{\vxi - \vxi'}_2$, con empates
    resueltos arbitrariamente. Entonces, los empates ocurren con probabilidad cero, y el hecho de que
    $S$ cubra $\Supp(\vx)$ a escala $\epsilon$ garantiza una distorsión no mayor
    que $\epsilon$; la tasa de este esquema es $\log_2 \cN_{\epsilon}(\Supp(\vx))$.

    Para la cota inferior, sea $0 < \delta < \tfrac{1}{2}$, y consideremos la
    variable aleatoria ensanchada $\vx_{\delta}$. Según el
    \Cref{lem:rate-distortion-lb-uos}, tenemos
    \begin{equation}
        R_{\delta + \epsilon}(\vx_{\delta})
        \leq
        R_{\epsilon}(\vx).
    \end{equation}
    Dado que $\vx_{\delta}$ tiene una densidad de Lebesgue uniforme, podemos entonces
    aplicar el límite inferior de Shannon, en la forma \eqref{eq:slb-appendix}, para obtener
    \begin{equation}\label{eq:slb-proof}
        \log_2 \volume(\Supp(\vx_{\delta}))
        - \log_2 \volume(B_{\delta + \epsilon})
        - O(\log D)
        \leq
        R_{\epsilon}(\vx).
    \end{equation}
    Por último, debemos hallar un límite inferior para la razón
    \begin{equation}
        \frac{
            \volume(\Supp(\vx_{\delta}))
        }
        {
            \volume(B_{\delta + \epsilon})
        }
    \end{equation}
    en términos del número de cobertura.
    Dado que $\Supp(\vx_{\delta}) = \Supp(\vx) + B_{\delta}$, donde $+$ denota aquí
    la suma de Minkowski, una aplicación estándar de los argumentos de límites de volumen
    (véase, por ejemplo, \cite[Proposition 4.2.12]{Vershynin2018-br}) da
    \begin{equation}
        \volume(\Supp(\vx_{\delta}))
        \geq
        \cN_{2\delta}(\Supp(\vx))
        \volume(B_{\delta}).
    \end{equation}
    Por lo tanto
    \begin{align}
        \frac{
            \volume(\Supp(\vx_{\delta}))
        }
        {
            \volume(B_{\delta + \epsilon})
        }
        &\geq
        \cN_{2\delta}(\Supp(\vx))
        \frac{
            \volume(B_{\delta})
        }
        {
            \volume(B_{\delta + \epsilon})
        }
        \\
        &=
        \cN_{2\delta}(\Supp(\vx))
        \left(
            \frac{
                \delta
            }
            {
                \delta + \epsilon
            }
        \right)^D.
    \end{align}
    Si se elige $\delta = \epsilon / 2$, se obtiene a partir del límite inferior de Shannon
    \eqref{eq:slb-proof} y las estimaciones anteriores
    \begin{equation}
        \log_2 \cN_{\epsilon}(\Supp(\vx))
        - O(D)
        \leq
        R_{\epsilon}(\vx),
    \end{equation}
    como se debía demostrar.

\end{proof}

\subsection{Demostración del Lema
\ref{lem:rate-distortion-lb-uos}}\label{sec:app-rate-dist-deferred-proofs}

\begin{proof}[(Demostración de \Cref{lem:rate-distortion-lb-uos})]
    Basta con demostrar que cualquier código para $\vx$ con una distorsión cuadrática esperada
    de $\epsilon^2$ produce un código para $\vx_{\delta}$ con la misma tasa y
    una distorsión ligeramente superior, si se elige un valor adecuado de $\delta$.
    Así pues, fijemos un código de este tipo para $\vx$, que alcance una tasa de $R$ y una distorsión cuadrática
    esperada de $\epsilon^2$. Denotamos por $\hat{\vx}$ la variable aleatoria
    reconstruida mediante este código, y por $\mathrm{q} : \Supp(\vx) \to \Supp(\vx)$
    la aplicación de codificación-decodificación asociada (es decir, $\hat{\vx}
    = \mathrm{q}(\vx)$).

    Sea ahora $\cS_k$ la hiperesfera de orden $k$ en el soporte de $\vx$.
    Existe una base ortonormal $\vU_k \in \R^{D \times d_k}$ tal que
    $\Span(\cS_k) = \Span(\vU_k)$.
    La siguiente descomposición ortogonal del conjunto de soporte $\cS$
    se utilizará repetidamente a lo largo de la demostración.
    Tenemos
    \begin{align}
        \cS_{\delta}
        &= \set{\vxi \in \R^D \given \exists k \in [K] \::\: \dist(\vxi,
        \cS_k) \leq \delta}
        \\
        &= \bigcup_{k \in [K]} \set{\vxi \in \R^D \given \dist(\vxi,
        \cS_k) \leq \delta}.
    \end{align}
    Mediante la proyección ortogonal, para cualquier $k \in [K]$, cualquier $\vxi \in \R^D$ puede
    escribirse como $\vxi = \vxi^{\|} + \vxi^{\perp}$, donde $\vxi^{\|} \in
    \Span(\cS_k)$
    y $\ip{\vxi^{\|}}{\vxi^\perp} = 0$.
    Entonces, para cualquier $\vxi' \in \cS_k$, tenemos
    \begin{align}
        \norm*{\vxi - \vxi'}_2^2
        =
        \norm*{\vxi^{\|} + \vxi^{\perp} - \vxi'}_2^2
        &=
        \norm*{\vxi^{\|}}_2^2 + \norm*{\vxi^{\perp}}_2^2 + \norm*{\vxi'}_2^2
        - 2 \ip*{\vxi^{\|}}{\vxi'}
        \\
        &\geq
        \norm*{\vxi^{\|}}_2^2 + \norm*{\vxi'}_2^2
        - 2 \ip*{\vxi^{\|}}{\vxi'}
        \\
        &=
        \norm*{\vxi^{\|} - \vxi'}_2^2.
    \end{align}
    Además, se sabe que para cualquier $\vxi^{\|}$ distinto de cero que pertenezca a $\Span(\cS_k)$,
    \begin{equation}
        \inf_{\vxi' \in \cS_k}\,
        \norm*{\vxi^{\|} - \vxi'}_2^2
        =
        \norm*{\vxi^{\|} - \frac{\vxi^{\|}}{\norm{\vxi^{\|}}_2}}_2^2.
    \end{equation}
    Si $\vxi^{\|}$ es cero, es evidente que la distancia anterior es igual a $1$
    para todo $\vxi' \in \cS_k$.
    Por lo tanto, si definimos un mapeo de proyección
    $\pi_{\cS_k}(\vxi)$
    mediante
    \begin{equation}
        \pi_{S_k}(\vxi) = \frac{\vU_k\vU_k^\top \vxi}{\norm{\vU_k^\top \vxi}_2}
    \end{equation}
    Para cualquier $\vxi \in \R^D$ tal que $\vU_k^\top \vxi \neq \mathbf{0}$, entonces
    $\pi_{\cS_k}(\vxi) = \argmin_{\vxi' \in \cS_k}\norm*{\vxi' - \vxi}_2$.
    Elegimos $0 < \delta < 1$, de modo que el conjunto engrosado $\cS_{\delta}$
    no contenga puntos $\vxi \in \R^D$ en los que alguna de los mapeos de proyección
    $\pi_{\cS_k}$ no esté bien definida.
    Por lo tanto, el conjunto engrosado $\cS_{\delta}$ satisface
    \begin{align}
        \cS_{\delta}
        &= \bigcup_{k \in [K]} \set*{\vxi \in \R^D \given
            \norm*{\vxi - \frac{\vU_k\vU_k^\top
            \vxi}{\norm{\vU_k^\top \vxi}_2}}_2
        \leq \delta}.
    \end{align}
    Estas distancias pueden reescribirse en términos de la descomposición ortogonal como
    \begin{align}
        \norm*{\vxi - \frac{\vU_k\vU_k^\top \vxi}{\norm{\vU_k^\top \vxi}_2}}_2^2
        &=
        \norm*{\vxi}_2^2 - 2 \norm{\vU_k^\top \vxi}_2 + 1
        \\
        &=
        \norm*{\vxi^{\|}}_2^2
        + \norm*{\vxi^{\perp}}_2^2
        - 2 \norm{\vU_k^\top \vxi^{\|}}_2 + 1
        \\
        &=
        \norm*{\vxi^{\perp}}_2^2
        + \left( \norm*{\vxi^{\|}}_2 - 1 \right)^2.
        \label{eq:distance-to-hypersphere}
    \end{align}
    A continuación demostraremos que cada uno de estos $\vxi \in \cS_{\delta}$ puede
    asociarse de forma única a una proyección sobre un único subespacio de la mezcla,
    lo que nos permitirá definir una proyección correspondiente sobre $\cS$.
    Dado un $\vxi \in \cS_{\delta}$, por lo anterior, podemos encontrar un subespacio
    $\vU_k$ tal que la descomposición ortogonal $\vxi = \vxi^{\|}_k
    + \vxi^{\perp}_k$ satisfaga
    \begin{equation}
        \norm*{\vxi^{\perp}_k}_2^2
        + \left( \norm*{\vxi^{\|}_k}_2 - 1 \right)^2
        \leq
        \delta^2.
    \end{equation}
    Consideremos la descomposición $\vxi = \vxi_j^{\|} + \vxi_j^{\perp}$ para algún $j
    \neq k$. Tenemos que
    \begin{align}
        \norm*{\vxi_j^{\|}}_2
        =
        \norm*{
            \vU_j \vU_j^\top \vxi
        }_2
        &=
        \norm*{
            \vU_j \vU_j^\top( \vU_k \vU_k^\top \vxi + (\vI - \vU_k \vU_k^\top)
            \vxi )
        }_2
        \\
        &=
        \norm*{
            \vU_j \vU_j^\top (\vI - \vU_k \vU_k^\top) \vxi
        }_2
        \\
        &\leq
        \norm*{
            (\vI - \vU_k \vU_k^\top) \vxi
        }_2
        =
        \norm*{
            \vxi_k^{\perp}
        }_2
        \leq \delta,
    \end{align}
    donde la segunda línea se basa en el supuesto de ortogonalidad de los subespacios
    $\vU_k$, y la tercera se basa en el hecho de que las proyecciones ortogonales son
    no expansivas.
    Por lo tanto, la $j$-ésima distancia cumple que
    \begin{equation}
        \norm*{\vxi^{\perp}_j}_2^2
        + \left( \norm*{\vxi^{\|}_j}_2 - 1 \right)^2
        \geq
        (1 - \delta)^2.
    \end{equation}
    Esto implica que si $0 < \delta < 1/2$, cada $\vxi \in \cS_{\delta}$ tiene
    un único subespacio más cercano en la mezcla.
    Por lo tanto, bajo esta condición, la siguiente aplicación $\pi_{\cS} : \cS_{\delta}
    \to \cS$ está bien definida:
    \begin{equation}
        \pi_{\cS}(\vxi)
        =
        \pi_{\cS_{k_\star}}(\vxi),
        \enspace\text{where}\enspace
        k_{\star} = \argmin_{k \in [K]}\,
        \dist(\vxi, \cS_{k}).
    \end{equation}

    Ahora, definimos un código para $\vx_{\delta}$ mediante
    \begin{equation}
        \hat{\vx}_{\delta} = \mathrm{q}( \pi_{\cS}(\vx_{\delta}) ).
    \end{equation}
    Es evidente que esto está relacionado con un código de tasa $R$ para $\vx_{\delta}$, ya que
    utiliza las correspondencias de codificación-decodificación del código de tasa $R$ para $\vx$. Tenemos
    que demostrar que presenta una distorsión reducida.
    Calculamos
    \begin{align}
        \bE\left[ \norm*{ \vx_{\delta} - \hat{\vx}_{\delta} }_2^2 \right]
        &=
        \bE\left[ \norm*{ \vx_{\delta} -
        \mathrm{q}(\pi_{\cS}(\vx_{\delta})) }_2^2 \right]
        \\
        &\leq
        \left(
            \bE\left[ \norm*{ \vx_{\delta} - \pi_{\cS}(\vx_{\delta}) }_2^2
            \right]^{1/2}
            + \bE\left[ \norm*{ \pi_{\cS}(\vx_{\delta})
            - \mathrm{q}(\pi_{\cS}(\vx_{\delta})) }_2^2 \right]^{1/2}
        \right)^2,\label{eq:distortion-decomposition}
    \end{align}
    donde la desigualdad utiliza la desigualdad de Minkowski.
    Ahora, por las \Cref{def:thickening-set,def:thickening-rv-uos}, tenemos
    determinísticamente
    \begin{equation}
        \norm*{ \vx_{\delta} - \pi_{\cS}(\vx_{\delta}) }_2^2
        \leq \delta^2,
    \end{equation}
    por lo que la esperanza también satisface esta estimación.
    Para el segundo término, bastará con caracterizar la densidad de la
    variable aleatoria $\pi_{\cS}(\vx_{\delta})$ como suficientemente cercana a la
    densidad de $\vx$ —la cual, tal y como implica el \Cref{assumption:union-of-spheres}, es
    una mezcla de distribuciones uniformes en cada subesfera $\cS_k$.
    Según el argumento anterior, cada punto $\vxi \in \cS_{\delta}$ puede asociarse
    a un único subespacio $\vU_k$, lo que significa que las componentes de la mezcla
    en la definición de $\cS_{\delta}$ (recordemos la
    \Cref{def:thickening-rv-uos}) no se superponen. Por lo tanto, la densidad
    $\pi_{\cS}(\vx_{\delta})$ puede caracterizarse estudiando el efecto de
    $\pi_{\cS_k}$ sobre la variable aleatoria condicional $\vx_{\delta}$, condicionada
    a ser extraída de $\cS_{k, \delta}$. Denotemos esta medida por $\mu_{k,
    \delta}$.
    Afirmamos que la proyección de esta medida bajo $\pi_{\cS_k}$
    es uniforme en $\cS_k$.
    Para ver que esto es cierto, recordamos la
    \Cref{eq:distance-to-hypersphere},
    que da la caracterización
    \begin{equation}
        \cS_{k, \delta} = \set*{
            \vxi^{\|} + \vxi^{\perp}
            \given
            \vxi^{\|} \in \Span(\vU_k),
            \vxi^{\perp} \in \Span(\vU_k)^\perp,
            \norm*{\vxi^{\perp}}_2^2
            + \left( \norm*{\vxi^{\|}}_2 - 1 \right)^2
            \leq
            \delta^{2}
        }.
    \end{equation}
    La distribución condicional en cuestión es uniforme en este conjunto; necesitamos
    demostrar que la proyección $\pi_{\cS_k}$ aplicada a esta variable aleatoria condicional
    da como resultado una variable aleatoria que es uniforme en $\cS_k$.
    Con respecto a estas coordenadas, hemos visto que $\pi_{\cS_k}(\vxi^\|
    + \vxi^\perp) = \vxi^\| / \norm{\vxi^{\|}}_2$.
    Por lo tanto, para cualquier $\vxi \in \cS_{k}$, tenemos que la preimagen de $\vxi$ en
    $\cS_{k, \delta}$ bajo $\pi_{\cS_k}$ es
    \begin{equation}
        \pi_{\cS_k}^{-1}(\vxi) = \set*{
            r\vxi + \vxi^\perp
            \given
            r > 0,
            \vxi^{\perp} \in \Span(\vU_k)^\perp,
            \norm*{\vxi^{\perp}}_2^2
            + \left( r - 1 \right)^2
            \leq
            \delta^{2}
        }.
        \label{eq:fiber-of-uos}
    \end{equation}
    Para demostrar que $(\pi_{\cS_k})_{\sharp} \mu_{k, \delta}$ es uniforme, necesitamos
    descomponer la integral de la densidad uniforme sobre $\cS_{k, \delta}$ de tal manera
    que quede claro que cada una de las fibras
    $\pi_{\cS_k}^{-1}(\vxi)$ (para cada $\vxi \in \cS_k$) «contribuye» por igual
    a la integral.\footnote{Más concretamente, esto equivale a descomponer
        la densidad uniforme en $\cS_{k, \delta}$ en una densidad condicional
        regular
        correspondiente a $\vxi \in \cS_{k}$, y demostrar que la densidad correspondiente
    en $\vxi$ es uniforme. La demostración deja claro que esto es cierto.}
    Por la \Cref{def:thickening-rv-uos} tenemos
    \begin{equation}
        \volume(\cS_{k, \delta})
        = \iint_{\Span(\vU_k) \times \Span(\vU_k)^\perp} \mathbf{1}_{
            \norm*{\vxi^{\perp}}_2^2
            + \left( \norm*{\vxi^{\|}}_2 - 1 \right)^2
            \leq
            \delta^{2}
        }
        \odif \vxi^{\|} \odif \vxi^{\perp}.
    \end{equation}
    En particular, la integración sobre las coordenadas ortogonales se factoriza.
    Sea $\odif \vtheta^{d}$ la medida uniforme (de Haar) sobre la esfera de
    radio $1$ en $\R^d$. Al convertir la integral $\vxi^{\|}$ a coordenadas polares,
    tenemos
    \begin{equation}
        \volume(\cS_{k, \delta})
        = \int_{[0, \infty)} \int_{\bS^{d_k-1}} \int_{\Span(\vU_k)^\perp}
        r^{d_k-1}
        \mathbf{1}_{
            \norm*{\vxi^{\perp}}_2^2
            + \left( r - 1 \right)^2
            \leq
            \delta^{2}
        }
        \odif r \odif \vtheta^{d_k} \odif \vxi^{\perp}.
    \end{equation}
    Si comparamos con la representación de fibra \eqref{eq:fiber-of-uos},
    vemos que necesitamos «integrar» sobre las componentes $r$ y $\vxi^\perp$
    de la integral anterior para verificar que el pushforward
    es uniforme.
    Pero esto es evidente, ya que la expresión anterior muestra que el valor de esta
    integral es independiente de $\vxi^\|$---o, de manera equivalente en este contexto, del
    valor de la componente esférica $\vtheta^{d_k}$.

    Por lo tanto, del argumento anterior se deduce que $\pi_{\cS}(\vx_{\delta})$ es
    uniforme.
    Dado que el supuesto sobre $\delta$ implica que las componentes de la mezcla en
    la distribución de $\vx_{\delta}$ no se superponen, los pesos de mezcla
    $\pi_k$ también se conservan en la imagen $\pi_{\cS}(\vx_{\delta})$, y en
    particular, la distribución de $\pi_{\cS}(\vx_{\delta})$ es igual a la
    distribución de $\vx$.
    Por lo tanto, el segundo término en la \Cref{eq:distortion-decomposition} satisface
    \begin{equation}
        \bE\left[ \norm*{ \pi_{\cS}(\vx_{\delta}) -
        \mathrm{q}(\pi_{\cS}(\vx_{\delta})) }_2^2 \right]
        =
        \bE\left[ \norm*{ \vx - \mathrm{q}(\vx) }_2^2 \right]
        \leq
        \epsilon^2,
    \end{equation}
    porque $\mathrm{q}$ es un código de distorsión $\epsilon$ para $\vx$.

    Así pues, hemos demostrado que el código hipotético de tasa $R$ y distorsión (esperada al cuadrado)
    $\epsilon^2$ para $\vx$ produce un código de tasa $R$ y distorsión (esperada al cuadrado)
    $\delta + \epsilon$ para $\vx_{\delta}$.
    Esto establece que
    \begin{equation}
        R_{\delta + \epsilon}(\vx_{\delta})
        \leq
        R_{\epsilon}(\vx),
    \end{equation}
    como se quería demostrar.

\end{proof}

\end{document}
